%Structure:title page; abstract; keywords; main text introduction, materials and methods, 
%results, discussion; acknowledgments; declaration of interest statement; references; appendices (as appropriate);
% table(s) and figure(s) and their captions placed within the text.
%Footnotes should be generally avoided, replaced with remarks in the text or placed in an appendix.
% When a footnote is absolutely necessary, please keep it brief and do not include mathematical expressions
%References within the text should give the authors' names, followed 
% by the year of publication in parentheses. References should be listed
% in alphabetical order at the end of the manuscript, after any appendixes
% but before any tables and figures. In this list, each reference should 
% contain, in the following order, the last name and initials of the author
% followed by those of the coauthors, year of publication, title of 
% reference, source, volume number, and page. Do not abbreviate journal
% titles. References to books must include the publisher's name and location.
% Only list references actually cited in the text.
\documentclass[12pt,letterpaper]{article}
\usepackage{tikz}
%\usetikzlibrary{datavisualization}
%\usetikzlibrary{datavisualization.formats.functions}
\usepackage[fleqn]{amsmath}
%\usepackage{amsrefs}
\usepackage{amssymb}% For Blackboard Fonts
\usepackage[indent=4ex,skip=.5\baselineskip]{parskip}
\usepackage{indentfirst}
\usepackage{centernot}
\usepackage{physics}
\usepackage{tabularx}
\usepackage{array}
\usepackage{url}
%\usepackage{bm}
\usepackage[perpage, symbol*]{footmisc}
\setfnsymbol{lamport}
\usepackage{cite}
%\usepackage{nicefrac}
% \usepackage{graphicx}
\usepackage{pgfplots}
\pgfplotsset{
    compat=newest,
}
\usepgfplotslibrary{groupplots}
\usepgfplotslibrary{external}
\tikzexternalize[prefix=RBAFig/]
%\usepackage[thinc]{esdiff}
%\usepackage{lipsum}
% Expectation symbol
\usepackage{amsfonts} %% <\m also included by amssymb
\DeclareMathSymbol{\m}{\mathbin}{AMSa}{"39}
\DeclareMathOperator*{\E}{\mathbb{E}}
\DeclareMathOperator{\Var}{Var}
\DeclareMathOperator*{\argmax}{argmax}
\DeclareMathOperator*{\argmin}{argmin}
\brokenpenalty=200\relax
\title{Robust Bayesian Statistics for Bernoulli Trials}
\author{Raleigh Bracewell Isadore Priour\\Raleigh@1337maps.com\\Tivy High School}
\date{April 2025}
\makeatletter
\renewcommand{\@makefnmark}{\hbox{\textsuperscript{\textbf{\tiny{\@thefnmark}}}}}
\makeatother
\usepackage{fancyhdr}
\usepackage{hyperref}
\hypersetup{
    colorlinks,
    citecolor=black,
    filecolor=black,
    linkcolor=black,
    urlcolor=black
}
\begin{document}
\pagestyle{fancy}
\fancyhf{}
\renewcommand{\headrulewidth}{0pt}
\fancyfoot[C]{\\\thepage}
\fancyfoot[L]{{\textsuperscript{\textsubscript{\copyright}}} 2025 Raleigh Priour All Rights Reserved}
\setlength{\mathindent}{0cm}
\setlength{\abovedisplayskip}{0pt}%
\setlength{\belowdisplayskip}{0pt}
\setlength{\abovedisplayshortskip}{0pt}%
\setlength{\belowdisplayshortskip}{0pt}
\allowdisplaybreaks
\maketitle
\section{Abstract}
Bayesian Analysis allows for a complete description of one's beliefs about the outcomes of events.
Although different people have different beliefs (priors), Robust Bayesian Analysis (RBA) allows one to start
from a baseline prior, for example an uninformative prior, and accepts all beliefs obtainable from that prior
given one saw at most the degree of prior certitude of beliefs, $d$, more points of evidence as reasonable. 
Inside the collection of reasonable probability distribution, nothing can be said about which prior/posterior is the \textit{correct} one. 
The point of Robust Bayesian Statistics is to convince ones audience of their result by encapsulating the viewer's belief's in their baseline.
Which departs from the usual goal of statistics to efficiently as possible determine conclusion from the data. 
In short, the goal of Robust Bayesian Statistics is to be convincing rather than efficent.
This paper goes through the derivation of the beta distribution as the conjugate prior for
the binomial distribution. Then it derives formulas and inequalities for Robust Bayesian Analysis of Bernoulli Trials.
In order to explain a graphing algorithm, which the author has implemented in python, 
for the maximum and minimum reasonable probabilities about the chance of success, $p_s$, of the binomial distribution. 
The author uses the last 30 years of the S\&P 500\cite{SP500} as an example to illustrate the methods covered in the paper.
\pagebreak
\tableofcontents
\par
This paper uses the nonstandard notation of $x!$ instead of the cumbersome $\Gamma(x+1)$ notation.
\section{Bayesian Inference for Bernoulli Trials}
\subsection{Introduction to the Beta Distribution}
\par
Consider that someone has just started out as a vacuum salesperson going door to door trying to 
figure out the chance\footnote{Probability is the correct term; however, this creates esoteric sentences 
like the probability that the probability is 20\%.} for someone opening the door and buying their product, $p_s$. 
They have never sold a vacuum before and assume the probability of the chances of success are equal likely; therefore,
their probability density function (pdf), $f$, for the chance of success is  $f(x)=1 \text{ for } x \in [0,1] $ where 
$x$ is the chance(probability) of success.
\par
Suppose the first two houses they went to did not buy a vacuum and let
$x=p_s$  therefore  $1\!\m\! x=1\!\m\! p_s=p_f$ where  $p_f$  is the probability of failure.
Given $p_f$ the probability of getting 2 failures and 0 successes is $(1\!\m\! p_f)^2$. \allowbreak The \mbox{(probability that $p_s=x$)$\propto (1\!\m\! x)^2$}. 
Integrating ones gets $ \displaystyle\int_0^1 (1\!\m\! x)^2 dx =\eval{ \m \frac{1}{3}(1 \m x)^3 + C}_0^1=\frac{1}{3} $ therefore 
\allowbreak $\displaystyle\Pr(p_s=x|\text{2 failures})\footnote{$\Pr(X|Y)$ means the probability of X given Y.}=3(1\!\m\! x)^2$ 
with the $3$ having no significance other than being the normalization constant. 
In general,\vspace{.7ex} \allowbreak \mbox{$\Pr(X | \text{Data})=\displaystyle\frac{\Pr(\text{Data} | X)\Pr(X)}{\Pr(\text{Data})}$}  (Baye's Theorem). 
\vspace{.3ex}This is the called the posterior distribution with $\Pr(p_s=x)=1$ being the prior distribution. 
Now assume one has $a$ successes and $b$ failures with the total number of observations, $n=a+b$. 
Then \allowbreak \mbox{$\Pr(p_s=x)\propto \Pr(\text{Success})^a \Pr(\text{Failure})^b=x^a (1\!\m\! x)^b$}. 
With \mbox{$\Pr(p_s=x)=\displaystyle\frac{x^a (1\!\m\! x)^b}{\int_{0}^{1} x^a (1\!\m\! x)^b dx}\quad\!$}\allowbreak let 
\mbox{$g(a,b)=\displaystyle\int_{0}^{1}x^a (1\!\m\! x)^b dx$}  \allowbreak
\begin{flalign*}
   &g(a,b)=\int_{0}^{1}x^a (1\!\m\! x)^{b}dx=\eval{\frac{x^{a+1}}{a+1}(1\!\m\! x)^b}_0^1 \!-\!  \int_{0}^{1}\frac{-b}{a+1}x^{a+1} (1\!\m\! x)^{b\m 1} dx  \\
   &\phantom{g(a,b)}=0+\frac{b}{a+1}g(a+1,b\!\m\! 1)\quad \text{for } b \ge 1  \\
   &\phantom{g(a,b)}=\frac{b}{a+1}\frac{b\!\m\! 1}{a+2}\cdot \ldots \cdot \frac{2}{a+b\!\m\! 1} \frac{1}{a+b}g(a+b,0) \\
   &\phantom{g(a,b)}=b!\frac{a!}{(a+b)!}\int_{0}^{1}x^{a+b}dx=\frac{a!b!}{(a+b)!}\frac{1}{a+b+1}=\frac{a!b!}{(a+b+1)!}\therefore\\
   &g(a,b)=\int_{0}^{1} x^a (1\!\m\! x)^b dx=\frac{a!b!}{(a+b+1)!}\\
   &\text{Therefore}\quad \frac{(a+b+1)!}{a!b!}\int_{0}^{1}x^a (1\!\m\! x)^b dx = \frac{(a+b+1)!}{a!b!}\frac{a!b!}{(a+b+1)!}=1\\
   &\text{So} \quad \Pr(p_s=x)=\frac{(a+b+1)!}{a!b!}x^a (1\!\m\! x)^b=\frac{(n+1)!}{a!b!}x^a (1\!\m\! x)^b\\ 
   &\text{From this we form the beta distribution } B(x;a,b)=\frac{(a+b+1)!}{a!b!}x^a (1\!\m\! x)^b\\ 
\end{flalign*}
\subsection{Descriptive Statistics of the Beta Distribution}
\begin{flalign*}
   &\E[B(a,b)]=\int_{0}^{1}\frac{(a+b+1)!}{a!b!}x^a (1\!\m\! x)^b x dx=\int_{0}^{1}\frac{(a+b+1)!}{a!b!}x^{a+1} (1\!\m\! x)^b dx\\
   &\phantom{\E[B(a,b)]}=\int_{0}^{1}\frac{(a+b+1)!}{a!b!}\frac{a+1+b+1}{a+1}\frac{a+1}{a+1+b+1}x^{a+1} (1\!\m\! x)^b dx\\
   &\phantom{\E[B(a,b)]}=\int_{0}^{1}\frac{(a+1+b+1)!}{(a+1)!b!}\frac{a+1}{a+b+2}x^{a+1} (1\!\m\! x)^b dx\\
   &\phantom{\E[B(a,b)]}=\frac{a+1}{a+b+2}\int_{0}^{1}\frac{(a+1+b+1)!}{(a+1)!b!}x^{a+1} (1\!\m\! x)^b dx\\
   &\E[B(a,b)]=\frac{a+1}{a+b+2}\cdot1=\frac{a+1}{a+b+2} \allowdisplaybreaks \\ 
   &\E[B(a,b)^2]=\int_{0}^{1}\frac{(a+b+1)!}{a!b!}x^a (1\!\m\! x)^b x^2 dx=\int_{0}^{1}\frac{(a+b+1)!}{a!b!}x^{a+2} (1\!\m\! x)^b dx\\   
   &\phantom{\E[B(a,b)]^2}=\int_{0}^{1}\frac{(a+b+1)!(a+b+2)(a+b+3)(a+1)(a+2)}{a!b!(a+1)(a+2)(a+b+2)(a+b+3)}x^{a+2} (1\!\m\! x)^b dx\\
   &\phantom{\E[B(a,b)]^2}=\int_{0}^{1}\frac{(a+b+3)!}{(a+2)!b!}\frac{(a+1)(a+2)}{(a+b+2)(a+b+3)}x^{a+2} (1\!\m\! x)^b dx\\
   &\phantom{\E[B(a,b)]^2}=\frac{(a+1)(a+2)}{(a+b+2)(a+b+3)}\int_{0}^{1}\frac{(a+b+3)!}{(a+2)!b!}x^{a+2} (1\!\m\! x)^b dx\\
   &\E[B(a,b)^2]=\frac{(a+1)(a+2)}{(a+b+2)(a+b+3)} \\
   &\Var[B(a,b)]=\E[B(a,b)^2]\!\m\! \E[B(a,b)]^2\\
   &\phantom{\Var[B(a,b)]}=\frac{(a+1)(a+2)}{(a+b+2)(a+b+3)}\!-\! \left(\frac{a+1}{a+b+2}\right)^2\\
   &\phantom{\Var[B(a,b)]}=\frac{a+1}{a+b+2}\left(\frac{a+2}{a+b+3}\!-\! \frac{a+1}{a+b+2} \right)\\
   &\phantom{\Var[B(a,b)]}=\frac{a+1}{a+b+2}\left(\frac{(a+2)(a+b+2)\!\m\! (a+1)(a+b+3)}{(a+b+3)(a+b+2)}\right)\\
   &\phantom{\Var[B(a,b)]}=\frac{a+1}{a+b+2}\left(\frac{b+1}{(a+b+3)(a+b+2)}\right)\\
   &\Var[B(a,b)]=\frac{(a+1)(b+1)}{(a+b+2)^2 (a+b+3)} \allowdisplaybreaks \\
   &\text{Mode for $a,b > 0$} \Leftrightarrow  \argmax_x \ln(B(x;a,b)) \Leftrightarrow \frac{\partial \ln(B(x;a,b))}{\partial x}=0\\
   &\ln(B(x;a,b))=\ln((a+b+1)!)\!\m\! \ln(a!)\!\m\! \ln(b!)+a\ln(x)+b\ln(x)\\
   &\frac{\partial \ln(B(x;a,b))}{\partial x}=\frac{a}{x}\!\m\! \frac{b}{1\!\m\! x}=0\\
   &\phantom{\frac{\partial \ln(B(x;a,b))}{\partial x}=}a(1\!\m\! x)\!\m\! bx=0 \therefore x=\frac{a}{a+b}\\
   &\text{Mode of $B(a,b)$}=\frac{a}{a+b} \quad \text{for $a,b > 0$}
\end{flalign*}
\subsection{Example}
\indent Suppose the salesperson went to 10 more houses and sold 2 vacuums, then their pdf for 
the chance (probability) of success is $B(x;2,10)$.
\noindent
\begin{center}
   \begin{tikzpicture}
      \begin{axis}[
         %title = {$B(x;2,10)$},  % whatever name you want
         xlabel = {Chance of Success},
         ylabel = {Probabilty},
         ymin = 0, ymax = 3.9,
         xmin=0, xmax=1.02,
         ytick distance=1,
         xtick distance=.2,
         width=.8\textwidth,
         height=.5\textwidth,
         axis x line = bottom,
         axis y line = left,
         clip=false,
         legend style={at={(1,.03)},anchor=south east,draw=none}
         ] 
         \addlegendimage{empty legend}
         \addlegendentry [text depth=]
             {$B(x;2,10)$}         
         \addplot[black,smooth,thick] table[x=x, y = y, col sep = comma] {B(x;2,10).csv};

      \end{axis}
  \end{tikzpicture}
\end{center}
\section{Framework for Robust Bayesian Statistics}
\indent There are many problems in which general consensus on the exact prior is not universal. For example, lets say 
one was modeling the height of Redwood trees as normally distributed with a prior mean, $\mu$, and standard deviation, $\sigma$.
Now suppose one person modeled $\Pr(\sigma=x) \propto \frac{1}{x}$ while a second person modeled standard deviation as 
$\Pr(\sigma=x)=x e^{\!\m  x}$.
Both of these are reasonable priors with the first one being scale invariant, while the second excludes implausible scales such as 
the tree's standard deviation being a mile. If data is collected and the posteriors for standard deviation disagree on parameters, 
such as the mean or credibility intervals, then said parameters cannot be derived from the data to much confidence since different 
reasonable starting points result in incompatible results. There is a popular phrase in statistics, 
\textit{let the data speak for itself}\footnote{This phrase is popular in frequentist statistics not Bayesian.},
by incorporating multiple reasonable priors and showing what these posteriors agree allows \textit{the data to speak for itself},
because it shows how one interprets the data does not significantly affect their conclusion. 
This is the definition I use for robust Bayesian statistics having statistical procedures that given a little more or less data 
does not significantly change the posterior.
\par
The problem this paper will be focused on is robust Bayesian statistics for the beta distribution.
We will determine the likelihood that  the S\&P 500 
will be higher at the end of the next trading day than it was at the last trading day %\footnote{without accounting for inflation.} 
based off the last 30 years the S\&P 500. The answer is surprisingly low between 52\%$-$55\%.
\par
For robust Bayesian statistics, define the reasonable set \allowbreak
\mbox{$\mathcal{R}(\theta)=\{\, \pi(\theta) \mid \pi \in \Pi \}\ $} for a parameter, $\theta$, with $\Pi$ being the set 
of all reasonable probability distribution. For many statistics 
\mbox{$\mathcal{R}(T)=[\inf \mathcal{R}(T),\sup \mathcal{R}(T)]$}\footnote{Interval may include or exclude its endpoints.}.  
However, the determination of any statistic, $T$, cannot go beyond the precision of \mbox{$\mathcal{R}(T)$}. 
This means robust Bayesian statistics is incapable of giving pointwise estimates. 
As annoying as this is, it is a good thing. 
Imagine asking five respected expert economic forecasters about their economic projections for the next 6 months. 
With three of the economist predicting moderately strong economic growth, one of them predicting slightly negative growth, 
and one of them predicting moderately strong negative economic growth, it is tempting to omit the moderately negative 
prediction and report that economic forecasters anticipate negligible to moderate economic growth.
However if years of school and experience in the field, producing respected economic predictions cannot lead one to the conclusion that the economy
will have a more positive outlook, then the data itself is not strong enough to suggest that conclusion. 
Bayesians generally do not like the phrase \textit{let the data speak for itself} because one's knowledge about the situation (one's prior)
changes what they take away from the data (their posterior).
\par
Now suppose one rotated an object with \mbox{3 sections: \textit{T,S,H}} 
and got \textit{S,T,H}. Now suppose it was spinner with sides of equal area labeled \textit{T,S,H}. 
This result would be unsurprising and would not give one knowledge 
that change their preconceived beliefs about the spinner, that each side is probably approximately equally likely as the others. 
However, if one took out a random nickel with \textit{H} for heads, \textit{T} for tails, and \textit{S} for the side/edge. 
One would be extremely surprised about getting \textit{S,T,H}, but would still believe that the coin 
will still most likely land head or tails and that observing the coin landing on its side was an extremely 
rare\footnote{However it is not as rare as one would think with a study predicting the probability to be about 1 in 6,000\cite{PhysRevE.48.2547}.} 
observation. It would be ridiculous to conclude from those 3 coins tosses that the probability of a nickel landing on its edge 
is about as likely as it landing on its head. It is only rational to incorporate one's knowledge about events in their model 
which can led the exact same data data leading to radically different conclusion as shown here.
Fortunately, robust Bayesian statistics does allow \textit{the data to speak for itself} because if all reasonable viewpoints (priors)
led to the a small $\mathcal{R}(\theta)$ set then one has usable clear conclusions which were not the results of someone's prior 
but the nature of the data. While it is still annoying to have a small range for a parameter instead of an exact number
such as the beliefs about the mean of $B(x;a,b)$ ranging from .68 to .72,
the real annoying truth is we want more precision than the data we collected can give us.
Going back to the vacuum salesperson example if the actual chance (probability) of success is 
$31.00\%$ it would take 32868 observations to be 95\% certain that the probability of successes is within $31\% \pm .5\%$ .
Nonetheless, if a person had a 100 experts in a room there would be very little they could unanimously agree 
on\footnote{Unless it was about theorems with mathematicians.}. Moreover, if someone just had a Bayesian and a frequentist they could 
not even agree on the correct definition of probability\footnote{It is the Bayesian one.}.
Therefore, it does matter on what priors one admits as reasonable to the problem with one aiming to incorporate the minimal range that includes 
all of the different arguable priors for some parameter. 
\section{RBA for Bernoulli Trials}
A good way to construct the collection of reasonable priors is to start with a baseline prior and then defining the reasonable priors as the
class of posteriors achievable based on seeing at most one's prior degree of certitude, $d$, additional observations. 
We shall now apply this analysis to Bernoulli trials.
Without loss of generality assume the baseline prior is $B(0,0)$ with one observing $a_s$ successes, 
$b_f$ failures, having prior degree of certitude, $d$, and \mbox{$\displaystyle{m=a_s\!+\!b_f\!+\!d}$}. 
If one wants the baseline prior to be $B(a^\prime,b^\prime)$ that is equivalent to observing $a_s+a^\prime$ successes 
and $b_f+b^\prime$ failures. 
\begin{flalign*}
   &\Pi=\{B(a_s+x,b_f+y) \mid x,y \in \mathbb{R}_{\ge0},x+y \le d\}
   &\intertext{Denote the ends of $\Pi$ as}
   &B_f(x)=B(x;a_s,b_f\!+\!d)=\frac{(m+1)!}{a_s!(b_f+d)!}x^{a_s}(1\!\m\! x)^{b_f+d} \text{ and}\\
   &B_s(x)=B(x;a_s\!+\!d,b_f)=\frac{(m+1)!}{(a_s+d)!b_f!}x^{a_s+d}(1\!\m\! x)^{b_f} \\
   &\mathcal{R}(\text{Mean})\left[\frac{a_s+1}{a_s+b_f+d+1},\frac{a_s+d+1}{a_s+d+b_f+1}\right] \\
   &\mathcal{R}(\text{Mode})=\left[\frac{a_s}{a_s+b_f+d},\frac{a_s+d}{a_s+d+b_f}\right]   
   &\intertext{Variance and standard deviation are tricker, because variance does not decrease monotonically with respect to $a$ or $b$ }
   &\Var[B(a,b)]=\frac{(a+1)(b+1)}{(a+b+2)^2(a+b+3)} \text{,  } \mu=\E[B(a,b)]=\frac{a+1}{a+b+2} \text{,  } n=a+b \\
   &\Var=\frac{\mu\cdot(1\!\m\! \mu)}{n+3} \implies d\!\phantom{.}\!\Var=\frac{1\!\m\! 2\mu}{n+3} d{\mu} \!\m\!  \frac{\mu(1\!\m\! \mu)}{(n+3)^2} d{n} \\
   &\mu=\frac{a+1}{a+b+2}=1\!\m\!\frac{b+1}{a+b+2} \implies \frac{\partial \mu}{\partial a}=\frac{b+1}{(a+b+2)^2} \text{ and } \frac{\partial \mu}{\partial b}=\m \frac{a+1}{(a+b+2)^2}\\
   &d{\mu}=\frac{1\!\m\! \mu}{n+2}da \!\m\!  \frac{\mu}{n+2}db \text{ and } dn=da+db \\
   &\frac{\partial \Var}{\partial a}=\frac{(1\!\m\! \mu)(1\!\m\! 2\mu)}{(n+3)(n+2)}\!\m\! \frac{\mu(1\!\m\! \mu)}{(n+3)^2}=\frac{1\!\m\! \mu}{n+3}\left(\frac{n+3\!\m\! (3n+8)\mu}{(n+2)(n+3)}\right)\\
   &\text{We want to find the values where an increase in $a$ does not decrease variance given $b$} \\
   &\frac{\partial \Var}{\partial a}=\frac{1\!\m\! \mu}{n+3}\left(\frac{n+3\!\m\! (3n+8)\mu}{(n+2)(n+3)}\right)\ge0 \\
   &n+3\!\m\! (3n+8)\mu\ge0 \\
   &\frac{n+3}{3n+8}\ge\mu \\
   &\frac{1}{3}+\frac{1}{9n+24}\ge1\!\m\! \frac{b+1}{n+2}\\
   &\frac{1}{9n+24}+\frac{b+1}{n+2}\ge\frac{2}{3}\\
   &(9n+24)b+10n+26\ge(6n+16)(n+2)\\
   &6n^2+(18\!\m\! 9b)n+6\!\m\! 24b \le0 \text{  carefully solving the inequality one gets}\\
   &n=a+b\le\sqrt{(\frac{3}{4}b+\frac{7}{6})^2 \!\m\! \frac{1}{9}}+\frac{3}{4}b\!\m\! \frac{3}{2}\\
   &a\le\sqrt{(\frac{3}{4}b+\frac{7}{6})^2 \!\m\! \frac{1}{9}}\!\m\! \frac{1}{4}b\!\m\! \frac{3}{2}\\
   &\text{For the case $a_s<b_f$ choose $d$ such that $a_s\!+\!d$ is the closest possible to} \\
   &\sqrt{(\frac{3}{4}b_f+\frac{7}{6})^2 \!\m\! \frac{1}{9}}\!\m\! \frac{1}{4}b_f\!\m\! \frac{3}{2}\\
   &\text{For the case $b_f<a_s$ choose $d$ such that $b_f\!+\!d$ is the closest possible to} \\
   &\sqrt{(\frac{3}{4}a_s+\frac{7}{6})^2 \!\m\! \frac{1}{9}}\!\m\! \frac{1}{4}a_s\!\m\! \frac{3}{2}\\
   &\text{For the case $a_s=b_f$ variance is already at it maximum with respect to $\mu$}
\end{flalign*}
\par
Credible intervals are still needed in robust Bayesian statistics. To determine $\mathcal{R(\text{Crl})}$ 
one just needs to find the beta distribution with the lowest lower bound, $b_l$, 
and the beta distribution with the highest upper bound, $b_u$, given the credible intervals, with \mbox{$\mathcal{R}(\text{Crl})=[b_l,b_u]$}. 
More formally, \\ \mbox{$\displaystyle{\mathcal{R}(\text{Crl}_\alpha(\theta))=\bigcup_{\pi\in\Pi}{\text{Crl}_\alpha(\theta;\pi)}}$}.
The Equal-Tailed Credible Interval (ETI) has the nice form, \allowbreak
\mbox{$\displaystyle{\mathcal{R}(\text{ETI}_{1\m \alpha}(p_s))=[Q_f(\alpha /2),Q_s(1-\alpha /2)]}$},
where $Q_f,Q_s$ are the quantile functions for $B_f,B_s$ respectively. 
To see this assume there exist a \mbox{$\displaystyle{B_t \in \Pi}$} with \mbox{$\displaystyle{B_t(x)=B(x;a_s\!+\!a_t,b_f\!+\!b_t)}$}
such that \mbox{$\displaystyle{Q_t(\alpha)<Q_f(\alpha)}$} for an $\alpha \in [0,1]$.
We are interested in the relationship\footnote{$\le$ and $\ge$.} between $B_f$ and $B_t$ which switches
only at points where $B_f(x)=B_t(x)$.
\begin{flalign*}
   &B(x;a_s,b_f\!+\!d)=B(x;a_s\!+\!a_t,b_f\!+\!b_t)\\
   &\frac{(a_s\!+\!b_f\!+\!d\!+\!1)!}{a_s!(b_f\!+\!d)!}x^{a_s}(1\!\m\!x)^{b_f+d}=\frac{(a_s\!+\!a_t\!+\!b_f\!+\!b_t\!+\!1)!}{(a_s\!+\!a_t)!(b_f\!+\!b_t)!}x^{a_s+a_t}(1\!\m\!x)^{b_f+b_t}\\
   &\frac{(a_s\!+\!b_f\!+\!d\!+\!1)!}{a_s!(b_f\!+\!d)!}(1\!\m\!x)^{d\m b_t}=\frac{(a_s\!+\!a_t\!+\!b_f\!+\!b_t\!+\!1)!}{(a_s\!+\!a_t)!(b_f\!+\!b_t)!}x^{a_t} \text{  for }x\in(0,1) \vspace{.3ex}
\end{flalign*}
\par
The left side of the equation is strictly decreasing and starts out greater than the right side, 
while the right side is strictly increasing and ends out greater than the left side. This implies
their exist an unique \mbox{$\displaystyle{x_t \in (0,1)}$} such that \mbox{$\displaystyle{B_f(x_t)=B_t(x_t)}$} therefore,
\begin{flalign*}
   &B_t(x) \le B_f(x) \text{  for  } x \in [0,x_t]\\
   &\implies F_t(x)\le F_f(x)\footnotemark{} \text{  for  }x \in [0,x_t] \text{ and}
\end{flalign*}
\footnotetext{$F_t$ and $F_f$ denote the cumulative distribution function of $B_t$ and $B_f$ respectively.}
\\[-1.8\baselineskip]
\begin{flalign*}
   &B_f(x) \le B_t(x) \text{  for  }x \in [x_t,1] \therefore\\
   &\int_{x}^{1}{B_f(y) dy} \le \int_{x}^{1}{B_t(y) dt}  \text{  for  }x \in [x_t,1] \\
   &\implies 1-F_f(x) \le 1-F_t(x) \text{  for  }x \in [x_t,1] \quad \therefore\\
   &F_t(x) \le F_f(x)\text{  for  }x \in [x_t,1] \text{ and for  }x \in [0,x_t] \quad \text{implying}\\
   &Q_f(\alpha)=F^{\m 1}_f(\alpha)\le F^{\m 1}_t(\alpha)=Q_t(\alpha) \text{  for  } \alpha \in [0,1]
\end{flalign*}mi
Using the same line of reasoning it can be shown that 
\begin{flalign*}
   &Q_t(\alpha)\le Q_s(\alpha) \text{  for  } \alpha \in [0,1]
\end{flalign*}
This also implies that \mbox{$\displaystyle{\mathcal{R}(\text{Median})=\left[Q_f\left(\frac{1}{2}\right),Q_s\left(\frac{1}{2}\right)\right]}$}
\section{Analytical Computational Methods\\ for Graphing}
Unlike normal Bayesian statistics in which one plots their posterior, we plot the regions of values 
$\pi(\theta)\footnote{$\pi$ denotes the probability distribution of the parameter, $\theta$.}$
can take. Define the upper posterior as \mbox{$\displaystyle{\pi_u (x)=\sup \mathcal{R}(\Pr(p_s=x))}$} and the lower posterior 
\mbox{$\displaystyle{\pi_l (x)=\inf \mathcal{R}(\Pr(p_s=x))}$}. 
$\pi_u$ starts as $B_f$ and stays $B_f$ until $x$ reaches the value $x_f$ 
where  ${{\frac{\partial B_f}{\partial a}}\big\rvert_{x_f,m}=0}$ from there $\pi_u$ continuously transition 
throughout the $a$ values ending at the point $x_s$ where $\eval{\frac{\partial B_s}{\partial a}}_{x_s,m}=0$.
Another way to put this is, given $x,n$ we are finding the $a$ value that maximizes $B(x;a,m\!\m\! a)$ given $a\in [a_s,a_s\!+\!d\,] $ 
with the extrema being outside $[a_s,a_s\!+\!d\,]$ for $x \notin [x_f,x_s]$ and 
the endpoints $B_f$ and $B_s$ being the optimal values until $x \in [x_f,x_s]$ in which an local extrema with respect to $a$ 
is in the interval $[a_s,a_s+d]$ with $a_s$ being an extrema at ${\frac{\partial B_f}{\partial a}}\big\rvert_{x,m}=0$ 
and $a_s+d$ being an extrema at $\eval{\frac{\partial B_s}{\partial a}}_{x,m}=0$.
However to differentiate $B_f,B_s$ with respect to $\mathit{a}$ we need to find the derivative of $x!$. 
To do this we first derive the extension of factorial function into $\mathbb{R}_{\ge0}.$  \vspace{.3ex}
\begin{flalign*}
   &\ln((M+n)!)=\ln(1\cdot2\cdot \ldots \cdot(M-1)\cdot M \cdot (M+1) \cdot (M+2) \cdot \ldots \cdot (M+n))\\
   &\phantom{\ln((M+n)!)}\,=\ln(M!)+ \ln((M+1) \cdot (M+2) \cdot \ldots \cdot (M+n))\\
   &\phantom{\ln((M+n)!)}\,=\ln(M!)+\ln(M^{n}+\frac{n(n\!+\!1)}{2}M^{n\m1}+\ldots+n!)\\
   &\phantom{\ln((M+n)!)}\,=\ln(M!)+\ln(M^{n}\cdot\left(1+\frac{n(n\!+\!1)}{2M}+\ldots+\frac{n!}{M^n}\right))\\[-.2em]
   &\ln((M+n)!)=\ln(M!)+n\ln(M)+\ln(1+\frac{n(n\!+\!1)}{2M}+\ldots+\frac{n!}{M^n})\\[-.7em]
   &\lim_{M \to \infty}\ln(1+\frac{n(n\!+\!1)}{2M}+\ldots+\frac{n!}{M^n})=0 \quad \text{therefore}\\
   &\lim_{M \to \infty} \ln((M+n)!) \to \ln(M!)+n\cdot\ln(M) \text{ we now assume this implies that}\\
   &\lim_{M \to \infty} \ln((M+x)!) \to \ln(M!)+x\cdot\ln(M)\\[-.2em]
   &\ln(x!)=\ln((M\!+\!x)!)\!\m\!\big(\!\ln((M\!+\!x)!)\!\m\! \ln(x!)\big)\\
   &\ln(x!)=\ln((M\!+\!x)!)\!\m\!\big(\!\ln(M\!+\!x)\!+\!\ln(M\!+\!x\!\m\!1)\!+\ldots+\!\ln(x\!+\!1)\!+\!\ln(x!)\!\m\! \ln(x!) \big)\\[-.2em]
   &\ln(x!)=\ln((M\!+\!x)!) - \sum_{k=1}^{M}{\ln(k\!+\!x)} \quad \therefore\\[-.5em]
   &\lim_{M \to \infty} \ln(x!)=\lim_{M \to \infty} \left( \ln((M\!+\!x)!) - \sum_{k=1}^{M}{\ln(k\!+\!x)}\right)\\
   &\ln(x!)=\lim_{M \to \infty} \left(\ln(M!)+x\ln(M) -\!\!\sum_{k=1}^{M}{\ln(k+x)}\right) \quad \therefore \\
   &\frac{d{\hspace{.1ex}\ln(x!)}}{d{x}}=\lim_{M \to \infty} \left(\ln(M) - \sum_{k=1}^{M}{\!\frac{1}{k+x}}\right)
   =\lim_{M \to \infty} \big(\ln(M)\!\m\! H_{M+x}\!+H_{x}\big) \quad \therefore \\
   &\frac{d{\hspace{.1ex}\ln(x!)}}{d{x}}=H_{x}\!\m\! \gamma\footnotemark{} \quad \text{\addtocounter{footnote}{-1} and} \quad \frac{d{x!}}{d{x}}=(H_x \!\m\!  \gamma)x! \\
\end{flalign*}
\stepcounter{footnote}\footnotetext{This is the Euler-Mascheroni Constant $\displaystyle{\gamma}=\lim_{x \to \infty}\left(H_x-\ln(x)\right)=0.5772156649\ldots$}
\\[-1.8\baselineskip]
To evaluate this at non-integer points we need to derive the continuation of the Harmonic Numbers.
\\[-1.5\baselineskip]
\begin{flalign*}
   &H_n\text{ denotes the nth Harmonic number with }H_n=\sum_{k=1}^{n}{\frac{1}{k}}=\frac{1}{n}+H_{n\m1} \\[-.35\baselineskip]
   &H_{M+n}=\sum_{k=1}^{M+n}{\frac{1}{k}}=\sum_{k=1}^{M}{\frac{1}{k}}+\sum_{k=1}^{n}{\frac{1}{M+k}}
   =H_M+\frac{n}{M}+\sum_{k=1}^{n}{\frac{1}{M+k}\!\m\!\frac{1}{M}}\\
   &\lim_{M \to \infty}\sum_{k=1}^{n}{\frac{1}{M+k}\!\m\!\frac{1}{M}}=0 \quad \therefore 
   \quad \lim_{M \to \infty} H_{M+n} \to H_M + \frac{n}{M}\\
   & \text{We assume this implies } \hspace{3.48em}\lim_{M \to \infty} H_{M+x} \to H_M + \frac{x}{M}\\[-.1em]
   &\text{We also assume that } H_{x+1}=\frac{1}{x\!+\!1}\!+\!H_{x}\Leftrightarrow H_{x+n}\!=H_{x}\!+\!\sum_{k=1}^{n}{\frac{1}{k\!+\!x}}\\[-.7em]
   &H_{x}=H_{M+x} \m (H_{M+x}\!\m\! H_x) = H_{M+x} \m \sum_{k=1}^{M}{\frac{1}{k\!+\!x}}\\[-.3em]
   &\lim_{M \to \infty}H_{x}=\lim_{M \to \infty}\left(H_{m+x} \m \sum_{k=1}^{M}{\frac{1}{k\!+\!x}} \right)=
   \lim_{M \to \infty} \left(H_{M} \!+\! \frac{x}{M} \m \sum_{k=1}^{M}{\frac{1}{x\!+\!k}}\right)\\
   &H_x=\lim_{M \to \infty}\left( \sum_{k=1}^{M}{\frac{1}{k}}+\frac{x}{M} \!\m\! \sum_{k=1}^{M}{\frac{1}{k+x}} \right)=\lim_{M \to \infty}\sum_{k=1}^{M}{\frac{1}{k}\!\m\! \frac{1}{k+x}}=\sum_{k=1}^{\infty}{\frac{1}{k}\!\m\! \frac{1}{k+x}}\\
   &\text{We are now ready to calculate } d{\hspace{.1ex}\ln(B)}\\
   &\frac{\partial \ln(B)}{\partial a}=\frac{\partial}{\partial a}\Big(\ln((a+b+1)!)\!\m\! \ln(a!)\!\m\! \ln(b!)+a\ln(x)+b\ln(1\!\m\! x)\Big)\\
   &\frac{\partial \ln(B)}{\partial a}=H_{a+b+1}\!\m\!\gamma\!\m\!\left(H_a\!\m\! \gamma\right)+\ln(x)\quad\text{Likewise}\\
   &\frac{\partial \ln(B)}{\partial b}=H_{a+b+1}\!\m\! H_b +\ln(1\!\m\!x) \therefore \\
   &d{\hspace{.1ex}\ln(B)}=\left(H_{a+b+1}\!\m\! H_a +\ln(x)\right)d{a}+\left(H_{a+b+1}\!\m\! H_b +\ln(1-x)\right)d{b} \\
   &\text{Now given the constraint }a+b=m\text{ using Lagrange multipliers to find }\\
   &\text{the constrained extrema of $\ln(B)$ one gets }\frac{\partial \ln(B)}{\partial a}=\lambda 
   \text{   and   } \frac{\partial \ln(B)}{\partial b}=\lambda \\
   &\text{therefore } \frac{\partial \ln(B)}{\partial a}=\frac{\partial \ln(B)}{\partial b} \text{ and we can replace }b\text{ with } m\!\m\!a\text{ getting the equation}\\
   &H_{m+1}\!\m\! H_a\!+\ln(x)=H_{m+1}\!\m\! H_{m\m a} \!+\ln(1-x) \Leftrightarrow H_{m\m a}\!\m\!H_a=\ln(\frac{1}{x}-1) \quad  \therefore\\
   &\ln(B(x;a,m\!\m\!a))\text{ is an extrema when }x=\frac{1}{e^{H_{m\m a}\m H_a}+1}%=\frac{e^{H_a\m H_{m\m a}}}{1+e^{H_a\m H_{m\m a}}}
   =1\!\m\! \frac{1}{1+e^{H_a\m H_{m\m a}}}\\
   &\text{Given these constraints $B_f$ is the maximum for }0\le x\le x_f\!=\!1\!\m\! \frac{1}{1+e^{H_{\small a_s}- H_{b_{\!{\small f}}+d}}}\\
   &\text{and $B_s$ is the maximum for }x_s\!=\!1\!\m\! \frac{1}{1+e^{H_{a_s+d} - H_{b_{\!f}}}}\le x\le 1\\
   &\text{For }x\in[x_f,x_s]\quad \pi_u(x(a))=B(x(a);a,m\!\m\! a) \quad \text{for }a\in[a_s,a_s+d]\\
   &\text{with } x(a)=1\!\m\! \frac{1}{1+e^{H_{a}\m H_{m \m a}}}\text{ and }\pi_u(x(a))=
   \frac{(m+1)!}{a!(m\!\m\! a)!}\frac{e^{a(H_a\m H_{m \m a})}}{(1+e^{H_a\m H_{m \m a}})^m}\\
   &\text{Given the constraint that $a+b=m$, the only extrema is a maximum; therefore, }\\
   &\pi_l(x)=\min(B_s(x),B_f(x))\text{ with }\pi_l \text{ starting at $B_s(x)$ and transitioning to $B_f(x)$ } \\
   &\text{at $x_c$ where }B_f(x_c)=B_s(x_c)\\
   &B_f(x_c)=\frac{(m+1)!} {a_s!(b_f+d)!}x_c^{a_s}(1\!\m\! x_c)^{b_f+d}=\frac{(m+1)!}{(a_s+d)!\,b_f!}x_c^{a_s+d}(1\!\m\! x_c)^{b_f}=B_s(x_c)\\
   &\Leftrightarrow \frac{(a_s+d)!}{a_s!}(1\!\m\! x_c)^d=\frac{(b_f+d)!}{b_f!}x_c^d \Leftrightarrow 1\!\m\! x_c=\sqrt[d]{\frac{(b_f+d)!a!}{b_f!(a_s+d)!}}x_c\\
   &\Leftrightarrow x_c=\frac{1}{1+\prod_{k=1}^{d}{\sqrt[d]{\frac{b_f+k}{a_s+k}}}} \quad \text{and}\quad x_c \in (x_f,x_s) 
\end{flalign*}
If $x_c \notin (x_f,x_s)$ then there would exist $x \in [0,1]$ where $\pi_u(x)=\pi_l(x)$\\\\
\allowbreak
\nopagebreak
\stepcounter{footnote}\footnotetext{The sections all use non-strict inequalities because at those points both definitions give the same answer.}
\begin{minipage}{\linewidth}
\hspace{-.5ex}From this one sees there are 4 sections\addtocounter{footnote}{-1}\footnotemark:
\begin{description}
   \item[$ \phantom{x_f}0 \le x \le x_f\phantom{1}$] $\pi_u=B_f$ and $\pi_l=B_s$
   \item[$ \phantom{0}x_f \le x \le x_c\phantom{1}$] $\pi_u(x(a))=B(x(a);a,m\!\m\! a)$ and $\pi_l=B_s$
   \item[$ \phantom{0}x_c \le x \le x_s\phantom{1}$] $\pi_u(x(a))=B(x(a);a,m\!\m\! a)$ and $\pi_l=B_f$
   \item[$ \phantom{0}x_s \le x \le 1\phantom{x_s}$] $\pi_u=B_s$ and $\pi_l=B_f$
\end{description}
\end{minipage}
\begin{subsection}{Example: The S\&P 500}
\par Looking at the last 30 years of the S\&P 500 in order to determine the likelihood of the S\&P 500 being higher the next trade day 
than the trade day before. I analyzed 2 cases using $d=100$ and $d=30$. For the case where $d=100$ what one can say about this probability 
is still true given one saw up to 100 additional days of S\&P 500 daily increases through decreases. 
For $d=30$ what one can say about this probability is still true given one saw up to 30 additional days of S\&P 500 daily 
increases through decreases. Lastly, the baseline prior is the initially belief that are probabilities are equally likely $B(x;0,0)$. 
Here is the robust Bayesian statistical analysis of the S\&P 500.\\[-1\baselineskip]
\begin{center}
\begin{tabularx}{\textwidth}{|>{\raggedleft\arraybackslash}X|>{\centering\arraybackslash}X|>{\centering\arraybackslash}X|}
   \hline
   S\&P 500&$d=30$ &$d=100$ \\ \hline
   Successes&$4063$&$4063$\\ \hline
   Failures&$3496$&$3496$\\ \hline
   Total&$7559$&$7559$ \\ \hline
   $\mathcal{R}(\text{Mean})$&$53.735\% \pm .395\%$&$53.701\% \pm 1.305\%$ \\ \hline
   $\mathcal{R}(\text{Mode})$&$53.736\% \pm .395\%$&$53.702\% \pm 1.306\%$ \\ \hline
   $\mathcal{R}(90\%\text{ ETI})$&$52.60\%-54.83\%$&$52.11\%-55.29\%$ \\ \hline
   $\mathcal{R}(95\%\text{ ETI})$&$52.41\%-55.05\%$&$51.93\%-55.47\%$ \\ \hline
   $\mathcal{R}(99\%\text{ ETI})$&$52.06\%-55.40\%$&$51.58\%-55.82\%$ \\ \hline
\end{tabularx}
\end{center}
\noindent
\\
\begin{flushleft}
   \begin{tikzpicture}
      \begin{groupplot}
         [group style={group size=2 by 2,horizontal sep=.07\textwidth, vertical sep=.1\textwidth,xlabels at=edge bottom,ylabels at=edge left
         },
         enlarge y limits=upper,enlarge y limits=0.04
         height=.27\textwidth,width=.40\textwidth,xlabel = {Probability of Being Higher},ylabel = {Probability},axis x line = bottom,axis y line=left
         ,legend style={at={(1,.06)},anchor=south east,draw=none}
         ,clip=false,clip mode=individual,scale only axis,legend style={at={(1,.06)},anchor=south east,draw=none}]
      \nextgroupplot[clip=false,title={$d=100$},xmax=1.02]
         \addplot[black,thick,tension=1] table[x=x, y = y, col sep = comma] {RBA 100 upper ends=True S&P 500.csv};
         \addplot[black,thick,loosely dashed,smooth] table[x=x, y = y, col sep = comma] {RBA 100 lower ends=True S&P 500.csv};
         \legend{$\pi_u$,$\pi_l$}
      \nextgroupplot[clip=false,title={$d=30$},xmax=1.02]
         \addplot[black,thick,tension=1] table[x=x, y = y, col sep = comma] {RBA 30 upper ends=True S&P 500.csv};
         \addplot[black,thick,loosely dashed,tension=1] table[x=x, y = y, col sep = comma] {RBA 30 lower ends=True S&P 500.csv}; 
         \legend{$\pi_u$,$\pi_l$}   
      \nextgroupplot[clip=false,title={$d=100$},legend style={at={(1.02,.7)}}]
         \addplot[black,thick,tension=1.5] table[x=x, y = y, col sep = comma] {RBA 100 upper ends=False S&P 500.csv};
         \addplot[black,thick,dashed,tension=1] table[x=x, y = y, col sep = comma] {RBA 100 lower ends=False S&P 500.csv};
         \legend{$\pi_u$,$\pi_l$}  
      \nextgroupplot[clip=false,title={$d=30$},legend style={at={(1,.7)}}]
         \addplot[black,thick,tension=1.5] table[x=x, y = y, col sep = comma] {RBA 30 upper ends=False S&P 500.csv};
         \addplot[black,thick,dashed,tension=1] table[x=x, y = y, col sep = comma] {RBA 30 lower ends=False S&P 500.csv};    
         \legend{$\pi_u$,$\pi_l$}            
   \end{groupplot}
  \end{tikzpicture}
\end{flushleft}
\par
From this even with $\pm100$ days of increases through decreases the probability of the S\&P 500 being higher the next trading day than the current 
trading day is between 51.6\% and 55.8\%. Showing that in the long-run the S\&P 500 is 
only slightly biased towards increasing from day to day.
\end{subsection}
\par The graph also show that a naive plotting algorithm that divides the 4 sections of the graph into equal spaced intervals 
will result in extremely bad graphs, because for the intervals outside $[x_f,x_s]$ the algorithm will sample  1 or 2 
points with non-negligible probability density, and the rest of the points it samples will have almost zero probability density.
\par I implemented robust Bayesian analysis of Bernoulli trials in python using matplotlib; however, the algorithm can be effectively 
implemented in any program or software capable of scientific computing and creating graphs. 
\par
If someone was plotting the curve $y=2x+2$ for $x\in [0,2]$ they would only need to provide the points:$(0,2),(2,4)$ to their graphing software 
as the line drawn between these two points would be the linear function that they wanted plotted. 
Therefore, what matters when plotting a function is not the function but the difference between the function and
the linear interpolation between the points sampled. The following algorithm calculates the next $x$ point where the area 
between the linear and quadratic approximation reaches some threshold, $\varepsilon$, and sample from there. 
Mathematically, $\displaystyle{f(x_0+x)\approx f(x_0)+f'(x_0)x+\frac{f''(x_0)}{2}x^2}$
and we want to find $\Delta x>0$ such that  $\displaystyle{\int_{0}^{\Delta x}{|f(x_0+t)\!\m\! (f(x_0)+f'(x_0)t)|dt} = \varepsilon}$ 
using the second order approximation we get $\displaystyle{\int_{0}^{\Delta x}{|f(x_0)+f'(x_0)t+\frac{f''(x_0)}{2}t^2\!\m\! (f(x_0)+f'(x_0)t)|dt} = \varepsilon}$  
so \allowbreak $\displaystyle{\int_{0}^{\Delta x}{\frac{|f''(x_0)|}{2}t^2dt} = \varepsilon}$  or 
$\displaystyle{\frac{|f''(x_0)|}{6}{\Delta x}^3=\varepsilon}$ 
so $\displaystyle{\Delta x=\sqrt[3]{\frac{6\varepsilon}{|f''(x_0)|}}}$. \\
Given $f$, $\Delta x_\text{avg} \propto \varepsilon^{1/3}$ therefore,
$\varepsilon^{\m 1/3}\propto \frac{1}{\Delta x_\text{avg}} \propto n_s$ where $n_s$ is the number of points sampled
Using regression from different values of $d$, $n$, and $p_s$, 
the asymptotical number of points sampled for $\pi_u$, $n_a$,  is approximately 
$\sqrt[3]{6.61889/\varepsilon}$ giving us $\varepsilon=6.61889/n_a^3$. 
The last thing we have to do is calculate $\displaystyle{f''}$ which comes in two variants $x$ as the independent variable
for $x\notin[x_f,x_s]$ and where $x$ is a function of $a$ for $x \in[x_f,x_s]$. 
We use this algorithm to sample $x$ points for $\pi_u$ and $\pi_l$ for $x \in[x_f,x_s]$
However, from the S\&P 500 we see that as $n$ becomes large, $\pi_u$ for $x \in[x_f,x_s]$ becomes flat and so we sample $x$ individually for 
$\pi_u$ and $\pi_l$, because in this region where $\pi_u$ is flat $\pi_l$ is the most interesting.
In order to implement said algorithm, we will derive $\displaystyle{\frac{d^2{\pi_u}}{d{x}^2},\frac{d^2{\pi_l}}{d{x}^2}}$. \vspace{.34ex}
\begin{flalign*}
   &\frac{\partial B(x;a,b)}{\partial x}=\frac{\partial}{\partial x}\left(\frac{(a+b+1)!}{a!b!}x^a (1\!\m\! x)^b\right)\\
   &\phantom{\frac{\partial B(x;a,b)}{\partial x}}=\frac{(a+b+1)!}{a!b!}\frac{a}{x}x^a(1\!\m\! x)^b+\frac{(a+b+1)!}{a!b!}\left(\frac{\m b}{1\!\m\!x}\right)x^a(1\!\m\! x)^b\\
   &\frac{\partial B(x;a,b)}{\partial x}=\left(\frac{a}{x}\!\m\! \frac{b}{1\!\m\! x}\right)B(x;a,b)\\
   &\frac{\partial^2{B(x;a,b)}}{{\partial x}^2}=\frac{\partial}{\partial x}\left(\left(\frac{a}{x}\!\m\! \frac{b}{1\!\m\! x}\right)B(x;a,b)\right)\\
   &\phantom{\frac{\partial^2{B(x;a,b)}}{{\partial x}^2}}=\left(\!\m \frac{a}{x^2}\!\m\! \frac{b}{(1\!\m\! x)^2}\right)B(x;a,b)+\left(\frac{a}{x}\!\m\! \frac{b}{1\!\m\! x}\right)^{\!\!2}\! B(x;a,b)\\
   &\frac{\partial^2{B(x;a,b)}}{{\partial x}^2}=\left[\left(\frac{a}{x}\!\m\! \frac{b}{1\!\m\! x}\right)^{\!\!2}\!\!\m\! \left(\frac{a}{x^2}+\frac{b}{(1\!\m\! x)^2}\right)\right]B(x;a,b)\\
   &\text{For } x \in[x_f,x_s]\quad x(a)=\frac{e^{H_a \m H_{m \m a}}}{1+e^{H_{a}\m H_{m \m a}}} \text{ but to compute }\frac{d{x}}{d{a}} \\
   &\text{we first need to be able to compute }\frac{d{H_t}}{d{t}} \text{ recall } H_t=\sum_{k=1}^{\infty}{\frac{1}{k}\!\m\! \frac{1}{k+t}}\\
   &\frac{d{H_t}}{d{t}}=\sum_{k=1}^{\infty}{\frac{1}{(k+t)^2}}=\psi_1(t)\footnotemark \text{ In general, }
   \frac{d^{\hspace{.05ex}n}{\!H_t}}{{d{t}}^n}=\psi_n(t)\footnotemark=\sum_{k=1}^{\infty}{\frac{(\!\m\! 1)^{n+1}n!}{(k+t)^{n+1}}} \quad \therefore
   \text{\addtocounter{footnote}{-1}\addtocounter{footnote}{-1}}
\end{flalign*} 
\stepcounter{footnote}\footnotetext{This is the trigamma function.} 
\stepcounter{footnote}\footnotetext{Using nonstandard notation of $k$ starting at 1 and not at 0 for the polygamma functions,$\psi_n$.}
\begin{flalign*}
   &\frac{d{x}}{d{a}} =\frac{e^{H_a\m H_{m \m a}}}{(1+e^{H_a\m H_{m \m a}})^2}(\psi_1(a)+\psi_1(m\!\m\! a))=\frac{\psi_1(a)+\psi_1(m\!\m\! a)}{1+e^{H_a\m H_{m \m a}}}x\\
   &\frac{d^2{x}}{d{a}^2}=\frac{\psi_2(a)\!\m\! \psi_2(m\!\m\! a)}{1+e^{H_a\m H_{m \m a}}}x\!\m\! \frac{e^{H_a\m H_{m \m a}}(\psi_1(a)+\psi_1(m\!\m\! a))^2}{(1+e^{H_a\m H_{m \m a}})^2}x
   +\frac{(\psi_1(a)+\psi_1(m\!\m\! a))^2}{(1+e^{H_a\m H_{m \m a}})^2}x\\
   &1\!\m\! 2x=1\!\m\! \frac{2e^{H_a\m H_{m \m a}}}{1+e^{H_a\m H_{m \m a}}}=\frac{1\!\m\! e^{H_a\m H_{m \m a}}}{1+e^{H_a\m H_{m \m a}}} \quad \therefore \\
   &\frac{d^2{x}}{d{a}^2}=\frac{x}{1+e^{H_a\m H_{m \m a}}}\left(\psi_2(a)\!\m\! \psi_2(m\!\m\! a)+(1\!\m\! 2x)\big(\psi_1(a)+\psi_1(m\!\m\! a)\big)^2\right)\\
   &\text{For }x\in [x_f,x_s] \:\:\: \pi_u(a)=B(x(a);a;m\!\m\! a)=\frac{(m+1)!}{a!(m\!\m\! a)!}\frac{e^{a(H_a\!\m\! H_{m \m a})}}{(1+e^{H_a\m H_{m \m a}})^m}\\
   &\ln{\pi_u(a)}=\ln((m+1)!)\!\m\! \ln(a!)\!\m\! \ln((m\!\m\! a)!)+a(H_a\!\m\! H_{m \m a})\!\m\! m\ln(1+e^{H_a\m H_{m \m a}})\\
   &\frac{d{\ln{\pi_u}}}{d{a}}=\gamma\!\m\! H_a\!\m\! \gamma\!+\!H_{m \m a}\!+\!H_a\!\m\! H_{m \m a}\!+\!a(\psi_1(a)\!+\!\psi_1(m\!\m\! a))\!\m\! mx(a)\!\cdot\!(\psi_1(a)\!+\!\psi_1(m\!\m\! a))\\
   &\frac{d{\ln{\pi_u}}}{d{a}}=(a\!\m\! mx(a))(\psi_1(a)+\psi_1(m\!\m\! a))\\
   &\frac{d^2{\ln{\pi_u}}}{d{a}^2}=\frac{d{}}{da}\left((a\!\m\! mx(a))(\psi_1(a)+\psi_1(m\!\m\! a))\right)\\
   &\phantom{\frac{d^2{\ln{\pi_u}}}{d{a}^2}}=\Big(1\!\m\! m\frac{\psi_1(a)\!+\!\psi_1(m\!\m\! a)}{1\!+\!e^{H_a\m H_{m \m a}}}x(a)\Big)(\psi_1(a)\!+\!\psi_1(m\!\m\! a))\!+\!(a\!\m\! mx(a))(\psi_2(a)\!\m\! \psi_2(m\!\m\! a))\\
   &\text{Surprisingly } \frac{d^2{f}}{{d{x}}^2} \centernot\iff 1\text{/}{\frac{d^2{x}}{{d{f}}^2}} \quad \text{Proof: }f(x)=x^2 \quad \frac{d^2{f}}{d{x}^2}=2 \quad 1\text{/}\frac{d^2{x}}{d{f}^2}=- 4x^3\\ 
   &\frac{d^2{\pi_u}}{{d{x}}^2}=\frac{d{}}{d{x}}\left(\frac{d{\pi_u}}{d{x}}\right)=\frac{d{}}{d{a}}\left(\frac{d{\pi_u}}{d{x}}\right)\cdot\frac{d{a}}{d{x}}
   =\frac{d{}}{d{a}}\left(\frac{d{\pi_u}}{d{a}}\frac{d{a}}{d{x}}\right)\cdot\frac{d{a}}{d{x}}\\
   &\phantom{\frac{d^2{\pi_u}}{{d{x}}^2}}=\left(\frac{d^2{\pi_u}}{{d{a}}^2}\frac{d{a}}{d{x}}+\frac{d{\pi_u}}{d{a}}\cdot\frac{d{}}{d{a}}\left(\frac{d{a}}{d{x}}\right)\right)\cdot\frac{d{a}}{d{x}}\\
   &\phantom{\frac{d^2{\pi_u}}{{d{x}}^2}}=\frac{d^2{\pi_u}}{{d{a}}^2}\left(\frac{d{a}}{d{x}}\right)^2+\frac{d{\pi_u}}{d{a}}\frac{d{a}}{d{x}}\cdot\frac{d{}}{d{a}}\left(\frac{1}{\frac{d{x}}{d{a}}}\right)
   =\frac{d^2{\pi_u}}{{d{a}}^2}\left(\frac{d{a}}{d{x}}\right)^2+\frac{d{\pi_u}}{d{a}}\frac{d{a}}{d{x}}\cdot\left(\!\m\! \frac{1}{\left(\frac{d{x}}{d{a}}\right)^2}\frac{d^2{x}}{{d{a}}^2} \right)\\
   &\frac{d^2{\pi_u}}{{d{x}}^2}=\frac{d^2{\pi_u}}{{d{a}}^2}\left(\frac{d{a}}{d{x}}\right)^{\!\!2}\!\m\! \frac{d{\pi_u}}{d{a}}\left(\frac{d{a}}{d{x}}\right)^{\!\!3}\frac{d^2{x}}{{d{a}}^2} \\
   &\frac{d^2{\ln \pi_u}}{{d{a}}^2}=\frac{d{}}{d{a}}\left(\frac{\pi_u'}{\pi_u}\right)=\frac{\pi_u''}{\pi_u}\!\m\! \frac{(\pi_u')^2}{\pi_u^2}=\frac{\pi_u''}{\pi_u}\!\m\! \left(\frac{d{\ln \pi_u}}{d{a}}\right)^{\!\!2}\quad\therefore\\
   &\frac{d^2{\pi_u}}{{d{x}}^2}=\pi_u\left(\frac{d^2{\ln \pi_u}}{{d{a}}^2}+\left(\frac{d{\ln \pi_u}}{d{a}}\right)^{\!\!2}\right)\quad\therefore\\
   &\frac{d^2{\pi_u}}{{d{x}}^2}=\pi_u\left(\frac{d{a}}{d{x}}\right)^{\!\!2}\left(\frac{d^2{\ln \pi_u}}{{d{a}}^2}+\frac{d{\ln \pi_u}}{d{a}}\left(\frac{d{\ln \pi_u}}{d{a}}\!\m\! \frac{d{a}}{d{x}}\frac{d^2{x}}{{d{a}}^2}\right)\right)\text{  for }x\in[x_f,x_s]\\
\end{flalign*}
Now that we know all of the derivatives we are ready to describe the \mbox{implementation:}
\pagebreak 
\begin{center}
   \begin{minipage}[t]{.45\textwidth}
      \begin{align*}
         &\varepsilon \gets \frac{6.61889}{n_a^3}\\
         &h\gets \sqrt[3]{6\varepsilon}\\
         &x \gets x_f\\
         &\pi_u,\pi_l=B_f,B_s\\
         &\textbf{While }x\ge0\text{:}\\
         &\qquad \textbf{Sample }\pi_u,\pi_l  \vspace{.5ex}\\
         &\qquad x \gets x-\frac{h}{\sqrt[3]{\left\lvert\frac{d^2{\pi_u}}{d{x}^2}\left(x\right)\right\rvert}}\\
         &a\gets a_s\\
         &\pi_u(a)=\pi(x(a),a,m\m a)\\
         &\textbf{While }a\le a_s+d\text{:}\\
         &\qquad x\gets x(a,m\m a)\\
         &\qquad \textbf{Sample } \pi_u\\
         &\qquad a \gets a+\frac{d{a}}{d{x}}\frac{h}{\sqrt[3]{\left\lvert\frac{d^2{\pi_u}}{d{x}^2}\left(x\right)\right\rvert}}
      \end{align*}
      \vspace{0pt}
   \end{minipage}
   \hfill
   \begin{minipage}[t]{.45\textwidth}
   \begin{align*}
      &x\gets x_c \\
      &\textbf{While }x\ge x_f\text{:}\\
      &\qquad \textbf{Sample } \pi_l\\
      &\qquad x \gets x-\frac{h}{\sqrt[3]{\left\lvert\frac{d^2{\pi_l}}{d{x}^2}\left(x\right)\right\rvert}}\\   
      &x\gets x_c \\
      &\pi_u,\pi_l=B_s,B_f\\
      &\textbf{While }x\le x_s\text{:}\\
      &\qquad \textbf{Sample } \pi_l\\
      &\qquad x \gets x+\frac{h}{\sqrt[3]{\left\lvert\frac{d^2{\pi_l}}{d{x}^2}\left(x\right)\right\rvert}}\\    
      &\textbf{While }x\le1\text{:}\\
      &\qquad \textbf{Sample }\pi_u,\pi_l  \vspace{.5ex}\\
      &\qquad x \gets x+\frac{h}{\sqrt[3]{\left\lvert\frac{d^2{\pi_u}}{d{x}^2}\left(x\right)\right\rvert}}   
   \end{align*}
   \vspace{0pt}
   \end{minipage}
\end{center}
\par
The author's implementation of the code can be found at \\
\mbox{\textit{\href{https://github.com/RaleighPriour/Robust-Bayesian-Analysis}{github.com/RaleighPriour/Robust-Bayesian-Analysis}\cite{Raleigh}}}
\nopagebreak \par
The last thing we will show is that the highest probability density value of $\pi_u$ is in the distribution with its mode furthest from $\frac{1}{2}$ which is 
$B_f(\frac{a_s}{m})$ when $2a_s\le n$ and $B_s(\frac{a_s+d}{m})$ when $2a_s\ge n$. As well as, \mbox{Mode $B_f\le x_f$} and 
\mbox{$x_s\le\text{Mode }B_s$}. We shall do this by looking at the structure of the beta family.\\
\begin{flalign*}
   &\text{By definition } \max\{B(x;a,b) \mid x \in [0,1]\}=B(\text{Mode }B(a,b);a,b).\\
   &\text{Taking the gradient of the beta family with the parameterization of the }\\
   &\text{mode denoted as $\nu$ and $n$ one gets}\\
   &B(\nu,n\nu,n(1\m \nu))=\frac{(n+1)!}{(n\nu)!(n(1 \m \nu))!}\nu^{n\nu}(1\m \nu)^{n(1 \m \nu)} \quad \therefore\\
   &\frac{\partial \ln B}{\partial \nu}=n(H_{n\nu} \m \gamma) \m n(H_{n(1 \m \nu)} \m \gamma) + n +n\ln(\nu) \m n \m n\ln(1\m \nu)\\
   &\frac{\partial \ln B}{\partial \nu}=n(H_{n\nu} \!-\! H_{n(1\m\nu)})+n\ln(\frac{\nu}{1\m\nu}) \\
   &\text{Now notice that }\frac{\partial \ln B}{\partial \nu}\text{ is }0\text{ and odd}\footnotemark\text{ about the point }\nu=\frac{1}{2}\text{ and is increasing}
\end{flalign*}
\footnotetext{$\frac{\partial \ln B}{\partial \nu}(\nu=\frac{1}{2}\m t)=\m\frac{\partial \ln B}{\partial \nu}(\nu=\frac{1}{2}+t)$}
\noindent
\!\!\!Therefore with respect to $\nu$ the beta distribution achieves a minimum at $\nu=\frac{1}{2}$ and symmetrically increases from $\nu=\frac{1}{2}$. 
\begin{flalign*}
   &\text{Now let us look at the beta distribution's behavior with respect to }n\\
   &\frac{\partial \ln B}{\partial n}=H_{n+1}\m\gamma\m \nu(H_{n\nu}\m\gamma) \m (1\m \nu)(H_{n(1\m \nu)} \m\gamma) + \nu\ln(\nu)+(1\m\nu)\ln(1\m\nu)\\
   &\frac{\partial \ln B}{\partial n}=H_{n+1} \m \nu H_{n\nu} \m (1\m\nu)H_{n(1\m\nu)} + \nu\ln(\nu)+(1\m\nu)\ln(1\m\nu)\\
   &\text{It is obvious that }\frac{\partial \ln B}{\partial n}\text{ is non-decreasing with respect to }n\text{.} \\
   &\text{Evaluating }\frac{\partial \ln B}{\partial n}\text{ at }n=0 \text{ one gets }\frac{\partial \ln B}{\partial n}\Bigr\rvert_{n=0}\!\!=1+\nu\ln(\nu)+(1\m\nu)\ln(1\m\nu)\\
   &\text{Which has a minimum of }1\m\ln(2)\text{ at $\textstyle \nu=\frac{1}{2}$ which one can easily verify. }
\end{flalign*}\noindent
Therefore the probability density at any beta distribution's modal value increases with respect to $n$, 
which makes intuitive sense the more data one collects the more informed their posterior distribution should be.
From the gradient of the beta distribution we see that the highest probability density value of $\pi_u$ is in the distribution with its mode furthest from $\frac{1}{2}$ which is 
$B_f(\frac{a_s}{m})$ when $2a_s\le n$ and $B_s(\frac{a_s+d}{m})$ when $2a_s\ge n$
% \begin{flalign*}
%    &\max{\pi_u}=\max\{\max\{\pi(x)\mid\pi\in\Pi\}\mid x\in [0,1]\}\\
%    &\phantom{\max{\pi_u}}=\max\{\max\{\pi(x)\mid x \in [0,1]\}\mid \pi\in \Pi\}\\
%    &\max{\pi_u}=\max\left\{B\left(\frac{a}{m};a,m \m a\right) \bigg\vert \: a\in [a_s,a_s+d_{\,}]\right\}\\
%    &y(a)=B\!\left(\text{Mode }B(a,m\!\m\! a);a,m \!\m\! a\right)=\frac{(m+1)!}{a!(m\m a)!}\frac{a^a(m\m a)^{(m\m a)}}{m^m}\\
%    &\ln{y(a)}=\ln((m+1)!)\m\ln(a!)\m\ln((m\m a)!)+a\ln(a)+(m\m a)\ln(m\m a)\m m\ln(m)\\
%    &\text{Since $\ln$ is a strictly increasing function and $y(a)>0$ this implies}\\
%    &\max{\pi_u}=\max{\{\ln{y(a)}\mid a\in [a_s,a_s+d]\}}\text{ which occurs when }\frac{d{\ln y}}{d{a}}=0\\
%    &\frac{d{\ln y}}{d{a}}=\m H_a+\gamma+H_{m\m a}\m\gamma+\ln(a)+1\m\ln(m\m a)\m 1\\
%    &\frac{d{\ln y}}{d{a}}=H_{m\m a}\m\ln(m\m a)\m\left(H_{a}\m\ln(a)\right)\\
%    &\text{Let }f(x)=H_x\m \ln(x)\text{ so }\frac{d{\ln y}}{d{a}}=f(m\m a)\m f(a)\\
%    &\ln{y}\text{ achieves a local extrema when }f(a)=f(m\m a)\text{ This only happens when }\\
%    &a=m\!\m\!a \Leftrightarrow a=\frac{m}{2}\text{ because }f\text{ is a strictly decreasing function and therefore injective\footnotemark}
%    \text{\addtocounter{footnote}{-1}}
% \end{flalign*} 
% \stepcounter{footnote}\footnotetext{Suppose $f(a)=f(m\m a)$ since $f$ is injective this implies $a=m\m a$.} 
% \vspace{-1.2ex}
% \begin{flalign*}   
%    &\text{To see that }f\text{ is strictly decreasing we will show that }f'(x)<0\text{ for all }x>0\\
%    &f'(x)=\frac{d{H_x}}{d{x}}\m \frac{1}{x}=\sum_{k=1}^{\infty}{\frac{1}{(k+x)^2}}\m\frac{1}{x}<\int_{0}^{\infty}{\frac{1}{(k+x)^2}d{\boldsymbol{k}}}\m\frac{1}{x}\\
%    &\phantom{f'(x)}<\eval{\m \frac{1}{k+x}+C}_{k=0}^{k=\infty}\m\frac{1}{x}=\frac{1}{x}\m\frac{1}{x}=0\quad\therefore\\
%    &f'(x)<0\text{ so $f$ is a strictly decreasing and injective function}\\
%    &\text{for }a<\frac{m}{2}\qquad m\m a>a  \quad \text{so} \quad f(m\m a)<f(a)\\
%    &\text{Since }\frac{d{\ln y}}{d{a}}=f(m\m a)\m f(a)\qquad\text{This shows }\frac{d{\ln y}}{d{a}}\text{ is decreasing from }\\
%    &a=0\text{ to }a=\frac{m}{2}\text{ and similarly that }\frac{d{\ln y}}{d{a}}\text{ is increasing from $a=\frac{m}{2}$ to $a=m$}
% \end{flalign*} \par
% Therefore, $\ln{y(a)}$ achieves a global minimum at $a=\frac{m}{2}$ for $a\in [0,m]$ and symmetrically increase from $\frac{m}{2}$,
% \allowbreak\mbox{$\ln{y(\frac{m}{2}-a)}=\ln{y(\frac{m}{2}-a)}$}. Therefore, $\pi_u$'s maximum given $a\in [a_s,a_s+d]$ is 
% $B_f(\frac{a_s}{m})$ when $2a_s\le n$ and $B_s(\frac{a_s+d}{m})$ when $2a_s\ge n$ 
\par
Finally, we will show \mbox{Mode $B_f\le x_f$} and \mbox{$x_s\le\text{Mode }B_s$}. 
By showing the following inequalities:
\begin{flalign*}
   &\qquad\text{Mode }B(a,m\!\m\!a)<x(a)<\frac{1}{2} \quad \text{when  }a<\frac{m}{2}\\
   &\qquad\frac{1}{2}=x(a)=\text{Mode }B(a,m\!\m\!a) \quad \text{when  }a=\frac{m}{2}\\
   &\qquad\frac{1}{2}<x(a)<\text{Mode }B(a,m\!\m\!a) \quad \text{when  }a>\frac{m}{2}\\
   &\text{Recall }x(a)=1\!\m\! \frac{1}{1+e^{H_a\m H_{m\m a}}} \text{  and  } \ln(\frac{x(a)}{1\!\m\!x(a)})=H_a\!\m\! H_{m\m a}\\
   &H_a\!\m\! H_{m\m a}=\sum_{k=1}^{\infty}{\frac{1}{k}\!\m\!\frac{1}{k+a}}-\sum_{k=1}^{\infty}{\frac{1}{k}\!\m\!\frac{1}{k+m\!\m\!a}}
   =\sum_{k=1}^{\infty}{\frac{1}{k+m\!\m\! a}\m \frac{1}{k+a}}\\
   &\text{assume }a>\frac{m}{2}\text{ so the sum is positive and therefore }\\
   &\ln(\frac{x(a)}{1\m x(a)})=\sum_{k=1}^{\infty}{\frac{1}{k+m\!\m\! a}\m \frac{1}{k+a}}<\int_{0}^{\infty}{\frac{1}{k+m\m a}\m\frac{1}{k+a}\,d{\boldsymbol{k}}}\\
   &\ln(\frac{x(a)}{1\m x(a)})<\eval{\ln\left(\frac{k+m\m a}{k+a}\right)+C}_{k=0}^{k=\infty}=\ln\left(\frac{a}{m\m a}\right)\\
   &\frac{x(a)}{1\m x(a)}<\frac{a}{m\m a} \Leftrightarrow x(a)<\frac{a}{m\!\m\!a}\m \frac{a}{m\!\m\!a}x(a) \Leftrightarrow \frac{m}{m\!\m\!a}x(a)<\frac{a}{m\!\m\!a}\\
   &x(a)<\frac{a}{m}=\text{Mode }B(a,m-a)\\
   &\text{If }a=\frac{m}{2}\text{ then }x(a)=1\!\m\! \frac{1}{1\!+\!e^{H_{m/2}\m H_{m\m m/2}}}=1\!\m\!\frac{1}{2}=\frac{1}{2}\text{ and Mode }B(\frac{m}{2},m)=\frac{1}{2}\\
   &\text{Finally, if }a<\frac{m}{2}\text{ then }\ln\left(\frac{x(a)}{1\!\m\!x(a)}\right)=H_{a}\!\m\!H_{m\m a}=-(H_{(m\m a)}\!\m\!H_{m\m (m\m a)}) \,\therefore\\
   &\ln\left(\frac{x(a)}{1\!\m\!x(a)}\right)=-\ln\left(\frac{x(m\!\m\!a)}{1\!\m\!x(m\!\m\!a)}\right)>
   -\ln\left(\frac{(m\!\m\! a)}{m-(m\!\m\! a)}\right)=-\ln\left(\frac{m\!\m\! a}{a}\right)\\
   &\frac{x(a)}{1\!\m\!x(a)}>\frac{a}{m\!\m\!a}\implies \text{Mode }B(a,m\!\m\!a)=\frac{a}{m}<x(a)<\frac{1}{2} 
\end{flalign*}
\section{Discussion}
\subsection{Implementation}
As said previously my implementation can be found on \href{https://github.com/RaleighPriour/Robust-Bayesian-Analysis}{\textbf{github\footnote{https://github.com/RaleighPriour/Robust-Bayesian-Analysis}}}. 
The implementation I have written can handle up to $n\approx2^{63}\approx10^{22}$ after that due to matplotlib 
and other libraries use of 64 bit numbers the graphs outputted become unstable; 
however, the actually statistics work for up to $n\approx10^{150}$ after that python is unable to handle the numerical values.
I also tried the procedure of calculating the  next sample point where the difference between the linear and 
quadratic approximation reach some threshold. Along with a method where the angle between the points plotted 
reached some threshold from 180°. Both methods worked significantly worse than the original cumulative error method.
\subsection{Case where $a_s=0$ or $b_f=0$}
Unlike the cases where $\displaystyle \lim_{n \to \infty } a_s,b_f \!\to \infty$ where the beta distribution approaches the normal distribution.
In these cases, $\pi_u$ and $\pi_l$ approach the same 2 relative curves that are most certainly not the 
normal distribution. Here is an example with $d=4$
\noindent
\begin{center}
   %% Creator: Matplotlib, PGF backend
%%
%% To include the figure in your LaTeX document, write
%%   \input{<filename>.pgf}
%%
%% Make sure the required packages are loaded in your preamble
%%   \usepackage{pgf}
%%
%% Also ensure that all the required font packages are loaded; for instance,
%% the lmodern package is sometimes necessary when using math font.
%%   \usepackage{lmodern}
%%
%% Figures using additional raster images can only be included by \input if
%% they are in the same directory as the main LaTeX file. For loading figures
%% from other directories you can use the `import` package
%%   \usepackage{import}
%%
%% and then include the figures with
%%   \import{<path to file>}{<filename>.pgf}
%%
%% Matplotlib used the following preamble
%%   \def#1{#1}
%%   \everymath=\expandafter{\the\everymath\displaystyle}
%%   \IfFileExists{scrextend.sty}{
%%     \usepackage[fontsize=10.000000pt]{scrextend}
%%   }{
%%     \renewcommand{\normalsize}{\fontsize{10.000000}{12.000000}\selectfont}
%%     \normalsize
%%   }
%%   
%%   \ifdefined\pdftexversion\else  % non-pdftex case.
%%     \usepackage{fontspec}
%%     \setmainfont{DejaVuSerif.ttf}[Path=\detokenize{C:/Users/prioural000/AppData/Local/Programs/Python/Python312/Lib/site-packages/matplotlib/mpl-data/fonts/ttf/}]
%%     \setsansfont{DejaVuSans.ttf}[Path=\detokenize{C:/Users/prioural000/AppData/Local/Programs/Python/Python312/Lib/site-packages/matplotlib/mpl-data/fonts/ttf/}]
%%     \setmonofont{DejaVuSansMono.ttf}[Path=\detokenize{C:/Users/prioural000/AppData/Local/Programs/Python/Python312/Lib/site-packages/matplotlib/mpl-data/fonts/ttf/}]
%%   \fi
%%   \makeatletter\@ifpackageloaded{underscore}{}{\usepackage[strings]{underscore}}\makeatother
%%
\begingroup%
\makeatletter%
\begin{pgfpicture}%
\pgfpathrectangle{\pgfpointorigin}{\pgfqpoint{2.602160in}{3.141564in}}%
\pgfusepath{use as bounding box, clip}%
\begin{pgfscope}%
\pgfsetbuttcap%
\pgfsetmiterjoin%
\definecolor{currentfill}{rgb}{1.000000,1.000000,1.000000}%
\pgfsetfillcolor{currentfill}%
\pgfsetlinewidth{0.000000pt}%
\definecolor{currentstroke}{rgb}{1.000000,1.000000,1.000000}%
\pgfsetstrokecolor{currentstroke}%
\pgfsetdash{}{0pt}%
\pgfpathmoveto{\pgfqpoint{0.000000in}{0.000000in}}%
\pgfpathlineto{\pgfqpoint{2.602160in}{0.000000in}}%
\pgfpathlineto{\pgfqpoint{2.602160in}{3.141564in}}%
\pgfpathlineto{\pgfqpoint{0.000000in}{3.141564in}}%
\pgfpathlineto{\pgfqpoint{0.000000in}{0.000000in}}%
\pgfpathclose%
\pgfusepath{fill}%
\end{pgfscope}%
\begin{pgfscope}%
\pgfsetbuttcap%
\pgfsetmiterjoin%
\definecolor{currentfill}{rgb}{1.000000,1.000000,1.000000}%
\pgfsetfillcolor{currentfill}%
\pgfsetlinewidth{0.000000pt}%
\definecolor{currentstroke}{rgb}{0.000000,0.000000,0.000000}%
\pgfsetstrokecolor{currentstroke}%
\pgfsetstrokeopacity{0.000000}%
\pgfsetdash{}{0pt}%
\pgfpathmoveto{\pgfqpoint{0.564660in}{0.521603in}}%
\pgfpathlineto{\pgfqpoint{2.502160in}{0.521603in}}%
\pgfpathlineto{\pgfqpoint{2.502160in}{2.831603in}}%
\pgfpathlineto{\pgfqpoint{0.564660in}{2.831603in}}%
\pgfpathlineto{\pgfqpoint{0.564660in}{0.521603in}}%
\pgfpathclose%
\pgfusepath{fill}%
\end{pgfscope}%
\begin{pgfscope}%
\pgfsetbuttcap%
\pgfsetroundjoin%
\definecolor{currentfill}{rgb}{0.000000,0.000000,0.000000}%
\pgfsetfillcolor{currentfill}%
\pgfsetlinewidth{0.803000pt}%
\definecolor{currentstroke}{rgb}{0.000000,0.000000,0.000000}%
\pgfsetstrokecolor{currentstroke}%
\pgfsetdash{}{0pt}%
\pgfsys@defobject{currentmarker}{\pgfqpoint{0.000000in}{-0.048611in}}{\pgfqpoint{0.000000in}{0.000000in}}{%
\pgfpathmoveto{\pgfqpoint{0.000000in}{0.000000in}}%
\pgfpathlineto{\pgfqpoint{0.000000in}{-0.048611in}}%
\pgfusepath{stroke,fill}%
}%
\begin{pgfscope}%
\pgfsys@transformshift{0.648567in}{0.521603in}%
\pgfsys@useobject{currentmarker}{}%
\end{pgfscope}%
\end{pgfscope}%
\begin{pgfscope}%
\definecolor{textcolor}{rgb}{0.000000,0.000000,0.000000}%
\pgfsetstrokecolor{textcolor}%
\pgfsetfillcolor{textcolor}%
\pgftext[x=0.648567in,y=0.424381in,,top]{\color{textcolor}{\sffamily\fontsize{10.000000}{12.000000}\selectfont\catcode`\^=\active\def^{\ifmmode\sp\else\^{}\fi}\catcode`\%=\active\def%{\%}${0.0}$}}%
\end{pgfscope}%
\begin{pgfscope}%
\pgfsetbuttcap%
\pgfsetroundjoin%
\definecolor{currentfill}{rgb}{0.000000,0.000000,0.000000}%
\pgfsetfillcolor{currentfill}%
\pgfsetlinewidth{0.803000pt}%
\definecolor{currentstroke}{rgb}{0.000000,0.000000,0.000000}%
\pgfsetstrokecolor{currentstroke}%
\pgfsetdash{}{0pt}%
\pgfsys@defobject{currentmarker}{\pgfqpoint{0.000000in}{-0.048611in}}{\pgfqpoint{0.000000in}{0.000000in}}{%
\pgfpathmoveto{\pgfqpoint{0.000000in}{0.000000in}}%
\pgfpathlineto{\pgfqpoint{0.000000in}{-0.048611in}}%
\pgfusepath{stroke,fill}%
}%
\begin{pgfscope}%
\pgfsys@transformshift{1.380700in}{0.521603in}%
\pgfsys@useobject{currentmarker}{}%
\end{pgfscope}%
\end{pgfscope}%
\begin{pgfscope}%
\definecolor{textcolor}{rgb}{0.000000,0.000000,0.000000}%
\pgfsetstrokecolor{textcolor}%
\pgfsetfillcolor{textcolor}%
\pgftext[x=1.380700in,y=0.424381in,,top]{\color{textcolor}{\sffamily\fontsize{10.000000}{12.000000}\selectfont\catcode`\^=\active\def^{\ifmmode\sp\else\^{}\fi}\catcode`\%=\active\def%{\%}${0.5}$}}%
\end{pgfscope}%
\begin{pgfscope}%
\pgfsetbuttcap%
\pgfsetroundjoin%
\definecolor{currentfill}{rgb}{0.000000,0.000000,0.000000}%
\pgfsetfillcolor{currentfill}%
\pgfsetlinewidth{0.803000pt}%
\definecolor{currentstroke}{rgb}{0.000000,0.000000,0.000000}%
\pgfsetstrokecolor{currentstroke}%
\pgfsetdash{}{0pt}%
\pgfsys@defobject{currentmarker}{\pgfqpoint{0.000000in}{-0.048611in}}{\pgfqpoint{0.000000in}{0.000000in}}{%
\pgfpathmoveto{\pgfqpoint{0.000000in}{0.000000in}}%
\pgfpathlineto{\pgfqpoint{0.000000in}{-0.048611in}}%
\pgfusepath{stroke,fill}%
}%
\begin{pgfscope}%
\pgfsys@transformshift{2.112834in}{0.521603in}%
\pgfsys@useobject{currentmarker}{}%
\end{pgfscope}%
\end{pgfscope}%
\begin{pgfscope}%
\definecolor{textcolor}{rgb}{0.000000,0.000000,0.000000}%
\pgfsetstrokecolor{textcolor}%
\pgfsetfillcolor{textcolor}%
\pgftext[x=2.112834in,y=0.424381in,,top]{\color{textcolor}{\sffamily\fontsize{10.000000}{12.000000}\selectfont\catcode`\^=\active\def^{\ifmmode\sp\else\^{}\fi}\catcode`\%=\active\def%{\%}${1.0}$}}%
\end{pgfscope}%
\begin{pgfscope}%
\definecolor{textcolor}{rgb}{0.000000,0.000000,0.000000}%
\pgfsetstrokecolor{textcolor}%
\pgfsetfillcolor{textcolor}%
\pgftext[x=1.533410in,y=0.234413in,,top]{\color{textcolor}{\sffamily\fontsize{10.000000}{12.000000}\selectfont\catcode`\^=\active\def^{\ifmmode\sp\else\^{}\fi}\catcode`\%=\active\def%{\%}$p_s$}}%
\end{pgfscope}%
\begin{pgfscope}%
\definecolor{textcolor}{rgb}{0.000000,0.000000,0.000000}%
\pgfsetstrokecolor{textcolor}%
\pgfsetfillcolor{textcolor}%
\pgftext[x=2.502160in,y=0.248302in,right,top]{\color{textcolor}{\sffamily\fontsize{10.000000}{12.000000}\selectfont\catcode`\^=\active\def^{\ifmmode\sp\else\^{}\fi}\catcode`\%=\active\def%{\%}$\times{10^{\ensuremath{-}5}}{}$}}%
\end{pgfscope}%
\begin{pgfscope}%
\pgfsetbuttcap%
\pgfsetroundjoin%
\definecolor{currentfill}{rgb}{0.000000,0.000000,0.000000}%
\pgfsetfillcolor{currentfill}%
\pgfsetlinewidth{0.803000pt}%
\definecolor{currentstroke}{rgb}{0.000000,0.000000,0.000000}%
\pgfsetstrokecolor{currentstroke}%
\pgfsetdash{}{0pt}%
\pgfsys@defobject{currentmarker}{\pgfqpoint{-0.048611in}{0.000000in}}{\pgfqpoint{-0.000000in}{0.000000in}}{%
\pgfpathmoveto{\pgfqpoint{-0.000000in}{0.000000in}}%
\pgfpathlineto{\pgfqpoint{-0.048611in}{0.000000in}}%
\pgfusepath{stroke,fill}%
}%
\begin{pgfscope}%
\pgfsys@transformshift{0.564660in}{0.626603in}%
\pgfsys@useobject{currentmarker}{}%
\end{pgfscope}%
\end{pgfscope}%
\begin{pgfscope}%
\definecolor{textcolor}{rgb}{0.000000,0.000000,0.000000}%
\pgfsetstrokecolor{textcolor}%
\pgfsetfillcolor{textcolor}%
\pgftext[x=0.289968in, y=0.573842in, left, base]{\color{textcolor}{\sffamily\fontsize{10.000000}{12.000000}\selectfont\catcode`\^=\active\def^{\ifmmode\sp\else\^{}\fi}\catcode`\%=\active\def%{\%}${0.0}$}}%
\end{pgfscope}%
\begin{pgfscope}%
\pgfsetbuttcap%
\pgfsetroundjoin%
\definecolor{currentfill}{rgb}{0.000000,0.000000,0.000000}%
\pgfsetfillcolor{currentfill}%
\pgfsetlinewidth{0.803000pt}%
\definecolor{currentstroke}{rgb}{0.000000,0.000000,0.000000}%
\pgfsetstrokecolor{currentstroke}%
\pgfsetdash{}{0pt}%
\pgfsys@defobject{currentmarker}{\pgfqpoint{-0.048611in}{0.000000in}}{\pgfqpoint{-0.000000in}{0.000000in}}{%
\pgfpathmoveto{\pgfqpoint{-0.000000in}{0.000000in}}%
\pgfpathlineto{\pgfqpoint{-0.048611in}{0.000000in}}%
\pgfusepath{stroke,fill}%
}%
\begin{pgfscope}%
\pgfsys@transformshift{0.564660in}{1.058710in}%
\pgfsys@useobject{currentmarker}{}%
\end{pgfscope}%
\end{pgfscope}%
\begin{pgfscope}%
\definecolor{textcolor}{rgb}{0.000000,0.000000,0.000000}%
\pgfsetstrokecolor{textcolor}%
\pgfsetfillcolor{textcolor}%
\pgftext[x=0.289968in, y=1.005949in, left, base]{\color{textcolor}{\sffamily\fontsize{10.000000}{12.000000}\selectfont\catcode`\^=\active\def^{\ifmmode\sp\else\^{}\fi}\catcode`\%=\active\def%{\%}${0.2}$}}%
\end{pgfscope}%
\begin{pgfscope}%
\pgfsetbuttcap%
\pgfsetroundjoin%
\definecolor{currentfill}{rgb}{0.000000,0.000000,0.000000}%
\pgfsetfillcolor{currentfill}%
\pgfsetlinewidth{0.803000pt}%
\definecolor{currentstroke}{rgb}{0.000000,0.000000,0.000000}%
\pgfsetstrokecolor{currentstroke}%
\pgfsetdash{}{0pt}%
\pgfsys@defobject{currentmarker}{\pgfqpoint{-0.048611in}{0.000000in}}{\pgfqpoint{-0.000000in}{0.000000in}}{%
\pgfpathmoveto{\pgfqpoint{-0.000000in}{0.000000in}}%
\pgfpathlineto{\pgfqpoint{-0.048611in}{0.000000in}}%
\pgfusepath{stroke,fill}%
}%
\begin{pgfscope}%
\pgfsys@transformshift{0.564660in}{1.490817in}%
\pgfsys@useobject{currentmarker}{}%
\end{pgfscope}%
\end{pgfscope}%
\begin{pgfscope}%
\definecolor{textcolor}{rgb}{0.000000,0.000000,0.000000}%
\pgfsetstrokecolor{textcolor}%
\pgfsetfillcolor{textcolor}%
\pgftext[x=0.289968in, y=1.438055in, left, base]{\color{textcolor}{\sffamily\fontsize{10.000000}{12.000000}\selectfont\catcode`\^=\active\def^{\ifmmode\sp\else\^{}\fi}\catcode`\%=\active\def%{\%}${0.4}$}}%
\end{pgfscope}%
\begin{pgfscope}%
\pgfsetbuttcap%
\pgfsetroundjoin%
\definecolor{currentfill}{rgb}{0.000000,0.000000,0.000000}%
\pgfsetfillcolor{currentfill}%
\pgfsetlinewidth{0.803000pt}%
\definecolor{currentstroke}{rgb}{0.000000,0.000000,0.000000}%
\pgfsetstrokecolor{currentstroke}%
\pgfsetdash{}{0pt}%
\pgfsys@defobject{currentmarker}{\pgfqpoint{-0.048611in}{0.000000in}}{\pgfqpoint{-0.000000in}{0.000000in}}{%
\pgfpathmoveto{\pgfqpoint{-0.000000in}{0.000000in}}%
\pgfpathlineto{\pgfqpoint{-0.048611in}{0.000000in}}%
\pgfusepath{stroke,fill}%
}%
\begin{pgfscope}%
\pgfsys@transformshift{0.564660in}{1.922924in}%
\pgfsys@useobject{currentmarker}{}%
\end{pgfscope}%
\end{pgfscope}%
\begin{pgfscope}%
\definecolor{textcolor}{rgb}{0.000000,0.000000,0.000000}%
\pgfsetstrokecolor{textcolor}%
\pgfsetfillcolor{textcolor}%
\pgftext[x=0.289968in, y=1.870162in, left, base]{\color{textcolor}{\sffamily\fontsize{10.000000}{12.000000}\selectfont\catcode`\^=\active\def^{\ifmmode\sp\else\^{}\fi}\catcode`\%=\active\def%{\%}${0.6}$}}%
\end{pgfscope}%
\begin{pgfscope}%
\pgfsetbuttcap%
\pgfsetroundjoin%
\definecolor{currentfill}{rgb}{0.000000,0.000000,0.000000}%
\pgfsetfillcolor{currentfill}%
\pgfsetlinewidth{0.803000pt}%
\definecolor{currentstroke}{rgb}{0.000000,0.000000,0.000000}%
\pgfsetstrokecolor{currentstroke}%
\pgfsetdash{}{0pt}%
\pgfsys@defobject{currentmarker}{\pgfqpoint{-0.048611in}{0.000000in}}{\pgfqpoint{-0.000000in}{0.000000in}}{%
\pgfpathmoveto{\pgfqpoint{-0.000000in}{0.000000in}}%
\pgfpathlineto{\pgfqpoint{-0.048611in}{0.000000in}}%
\pgfusepath{stroke,fill}%
}%
\begin{pgfscope}%
\pgfsys@transformshift{0.564660in}{2.355031in}%
\pgfsys@useobject{currentmarker}{}%
\end{pgfscope}%
\end{pgfscope}%
\begin{pgfscope}%
\definecolor{textcolor}{rgb}{0.000000,0.000000,0.000000}%
\pgfsetstrokecolor{textcolor}%
\pgfsetfillcolor{textcolor}%
\pgftext[x=0.289968in, y=2.302269in, left, base]{\color{textcolor}{\sffamily\fontsize{10.000000}{12.000000}\selectfont\catcode`\^=\active\def^{\ifmmode\sp\else\^{}\fi}\catcode`\%=\active\def%{\%}${0.8}$}}%
\end{pgfscope}%
\begin{pgfscope}%
\pgfsetbuttcap%
\pgfsetroundjoin%
\definecolor{currentfill}{rgb}{0.000000,0.000000,0.000000}%
\pgfsetfillcolor{currentfill}%
\pgfsetlinewidth{0.803000pt}%
\definecolor{currentstroke}{rgb}{0.000000,0.000000,0.000000}%
\pgfsetstrokecolor{currentstroke}%
\pgfsetdash{}{0pt}%
\pgfsys@defobject{currentmarker}{\pgfqpoint{-0.048611in}{0.000000in}}{\pgfqpoint{-0.000000in}{0.000000in}}{%
\pgfpathmoveto{\pgfqpoint{-0.000000in}{0.000000in}}%
\pgfpathlineto{\pgfqpoint{-0.048611in}{0.000000in}}%
\pgfusepath{stroke,fill}%
}%
\begin{pgfscope}%
\pgfsys@transformshift{0.564660in}{2.787138in}%
\pgfsys@useobject{currentmarker}{}%
\end{pgfscope}%
\end{pgfscope}%
\begin{pgfscope}%
\definecolor{textcolor}{rgb}{0.000000,0.000000,0.000000}%
\pgfsetstrokecolor{textcolor}%
\pgfsetfillcolor{textcolor}%
\pgftext[x=0.289968in, y=2.734376in, left, base]{\color{textcolor}{\sffamily\fontsize{10.000000}{12.000000}\selectfont\catcode`\^=\active\def^{\ifmmode\sp\else\^{}\fi}\catcode`\%=\active\def%{\%}${1.0}$}}%
\end{pgfscope}%
\begin{pgfscope}%
\definecolor{textcolor}{rgb}{0.000000,0.000000,0.000000}%
\pgfsetstrokecolor{textcolor}%
\pgfsetfillcolor{textcolor}%
\pgftext[x=0.234413in,y=1.676603in,,bottom,rotate=90.000000]{Probabilty}
\end{pgfscope}%
\begin{pgfscope}%
\definecolor{textcolor}{rgb}{0.000000,0.000000,0.000000}%
\pgfsetstrokecolor{textcolor}%
\pgfsetfillcolor{textcolor}%
\pgftext[x=0.564660in,y=2.873270in,left,base]{\color{textcolor}{\sffamily\fontsize{10.000000}{12.000000}\selectfont\catcode`\^=\active\def^{\ifmmode\sp\else\^{}\fi}\catcode`\%=\active\def%{\%}$\times{10^{6}}{}$}}%
\end{pgfscope}%
\begin{pgfscope}%
\pgfpathrectangle{\pgfqpoint{0.564660in}{0.521603in}}{\pgfqpoint{1.937500in}{2.310000in}}%
\pgfusepath{clip}%
\pgfsetrectcap%
\pgfsetroundjoin%
\pgfsetlinewidth{1.505625pt}%
\definecolor{currentstroke}{rgb}{0.000000,0.000000,0.000000}%
\pgfsetstrokecolor{currentstroke}%
\pgfsetdash{}{0pt}%
\pgfpathmoveto{\pgfqpoint{0.652728in}{2.726603in}}%
\pgfpathlineto{\pgfqpoint{0.663055in}{2.583606in}}%
\pgfpathlineto{\pgfqpoint{0.673633in}{2.447214in}}%
\pgfpathlineto{\pgfqpoint{0.684476in}{2.317272in}}%
\pgfpathlineto{\pgfqpoint{0.695596in}{2.193627in}}%
\pgfpathlineto{\pgfqpoint{0.707009in}{2.076125in}}%
\pgfpathlineto{\pgfqpoint{0.718731in}{1.964609in}}%
\pgfpathlineto{\pgfqpoint{0.730779in}{1.858927in}}%
\pgfpathlineto{\pgfqpoint{0.730779in}{1.858927in}}%
\pgfpathlineto{\pgfqpoint{0.739245in}{1.793199in}}%
\pgfpathlineto{\pgfqpoint{0.748483in}{1.732061in}}%
\pgfpathlineto{\pgfqpoint{0.758551in}{1.675120in}}%
\pgfpathlineto{\pgfqpoint{0.769513in}{1.622020in}}%
\pgfpathlineto{\pgfqpoint{0.781434in}{1.572444in}}%
\pgfpathlineto{\pgfqpoint{0.794386in}{1.526102in}}%
\pgfpathlineto{\pgfqpoint{0.808441in}{1.482734in}}%
\pgfpathlineto{\pgfqpoint{0.823678in}{1.442103in}}%
\pgfpathlineto{\pgfqpoint{0.840178in}{1.403996in}}%
\pgfpathlineto{\pgfqpoint{0.858026in}{1.368220in}}%
\pgfpathlineto{\pgfqpoint{0.877314in}{1.334597in}}%
\pgfpathlineto{\pgfqpoint{0.898135in}{1.302968in}}%
\pgfpathlineto{\pgfqpoint{0.920588in}{1.273187in}}%
\pgfpathlineto{\pgfqpoint{0.944777in}{1.245120in}}%
\pgfpathlineto{\pgfqpoint{0.970812in}{1.218645in}}%
\pgfpathlineto{\pgfqpoint{0.998804in}{1.193652in}}%
\pgfpathlineto{\pgfqpoint{1.028872in}{1.170037in}}%
\pgfpathlineto{\pgfqpoint{1.061141in}{1.147708in}}%
\pgfpathlineto{\pgfqpoint{1.095740in}{1.126578in}}%
\pgfpathlineto{\pgfqpoint{1.132802in}{1.106567in}}%
\pgfpathlineto{\pgfqpoint{1.172470in}{1.087603in}}%
\pgfpathlineto{\pgfqpoint{1.214888in}{1.069619in}}%
\pgfpathlineto{\pgfqpoint{1.260211in}{1.052552in}}%
\pgfpathlineto{\pgfqpoint{1.308596in}{1.036344in}}%
\pgfpathlineto{\pgfqpoint{1.308829in}{1.036270in}}%
\pgfpathlineto{\pgfqpoint{1.339408in}{1.025061in}}%
\pgfpathlineto{\pgfqpoint{1.371978in}{1.010141in}}%
\pgfpathlineto{\pgfqpoint{1.407286in}{0.991248in}}%
\pgfpathlineto{\pgfqpoint{1.446638in}{0.967793in}}%
\pgfpathlineto{\pgfqpoint{1.492887in}{0.938265in}}%
\pgfpathlineto{\pgfqpoint{1.556801in}{0.896303in}}%
\pgfpathlineto{\pgfqpoint{1.627998in}{0.850897in}}%
\pgfpathlineto{\pgfqpoint{1.679052in}{0.820541in}}%
\pgfpathlineto{\pgfqpoint{1.726272in}{0.794655in}}%
\pgfpathlineto{\pgfqpoint{1.771947in}{0.771841in}}%
\pgfpathlineto{\pgfqpoint{1.817105in}{0.751516in}}%
\pgfpathlineto{\pgfqpoint{1.862374in}{0.733353in}}%
\pgfpathlineto{\pgfqpoint{1.908215in}{0.717135in}}%
\pgfpathlineto{\pgfqpoint{1.955016in}{0.702700in}}%
\pgfpathlineto{\pgfqpoint{2.003139in}{0.689913in}}%
\pgfpathlineto{\pgfqpoint{2.052949in}{0.678659in}}%
\pgfpathlineto{\pgfqpoint{2.104836in}{0.668831in}}%
\pgfpathlineto{\pgfqpoint{2.159239in}{0.660328in}}%
\pgfpathlineto{\pgfqpoint{2.216670in}{0.653054in}}%
\pgfpathlineto{\pgfqpoint{2.277746in}{0.646911in}}%
\pgfpathlineto{\pgfqpoint{2.343228in}{0.641805in}}%
\pgfpathlineto{\pgfqpoint{2.414092in}{0.637641in}}%
\pgfusepath{stroke}%
\end{pgfscope}%
\begin{pgfscope}%
\pgfpathrectangle{\pgfqpoint{0.564660in}{0.521603in}}{\pgfqpoint{1.937500in}{2.310000in}}%
\pgfusepath{clip}%
\pgfsetbuttcap%
\pgfsetroundjoin%
\pgfsetlinewidth{1.505625pt}%
\definecolor{currentstroke}{rgb}{0.000000,0.000000,0.000000}%
\pgfsetstrokecolor{currentstroke}%
\pgfsetdash{{5.550000pt}{2.400000pt}}{0.000000pt}%
\pgfpathmoveto{\pgfqpoint{0.652728in}{0.626603in}}%
\pgfpathlineto{\pgfqpoint{0.663055in}{0.626611in}}%
\pgfpathlineto{\pgfqpoint{0.673633in}{0.626668in}}%
\pgfpathlineto{\pgfqpoint{0.684476in}{0.626858in}}%
\pgfpathlineto{\pgfqpoint{0.695596in}{0.627298in}}%
\pgfpathlineto{\pgfqpoint{0.707009in}{0.628136in}}%
\pgfpathlineto{\pgfqpoint{0.718731in}{0.629543in}}%
\pgfpathlineto{\pgfqpoint{0.730779in}{0.631706in}}%
\pgfpathlineto{\pgfqpoint{0.748017in}{0.636316in}}%
\pgfpathlineto{\pgfqpoint{0.771607in}{0.645974in}}%
\pgfpathlineto{\pgfqpoint{0.795131in}{0.659815in}}%
\pgfpathlineto{\pgfqpoint{0.819239in}{0.678402in}}%
\pgfpathlineto{\pgfqpoint{0.844721in}{0.702545in}}%
\pgfpathlineto{\pgfqpoint{0.872995in}{0.733890in}}%
\pgfpathlineto{\pgfqpoint{0.934558in}{0.812405in}}%
\pgfpathlineto{\pgfqpoint{0.972661in}{0.862820in}}%
\pgfpathlineto{\pgfqpoint{0.972661in}{0.862820in}}%
\pgfpathlineto{\pgfqpoint{0.993556in}{0.831407in}}%
\pgfpathlineto{\pgfqpoint{1.015469in}{0.802940in}}%
\pgfpathlineto{\pgfqpoint{1.038503in}{0.777273in}}%
\pgfpathlineto{\pgfqpoint{1.062777in}{0.754256in}}%
\pgfpathlineto{\pgfqpoint{1.088429in}{0.733742in}}%
\pgfpathlineto{\pgfqpoint{1.115625in}{0.715581in}}%
\pgfpathlineto{\pgfqpoint{1.144558in}{0.699628in}}%
\pgfpathlineto{\pgfqpoint{1.175460in}{0.685734in}}%
\pgfpathlineto{\pgfqpoint{1.208615in}{0.673753in}}%
\pgfpathlineto{\pgfqpoint{1.244368in}{0.663538in}}%
\pgfpathlineto{\pgfqpoint{1.283154in}{0.654943in}}%
\pgfpathlineto{\pgfqpoint{1.308829in}{0.650385in}}%
\pgfpathlineto{\pgfqpoint{1.339408in}{0.645903in}}%
\pgfpathlineto{\pgfqpoint{1.371978in}{0.642054in}}%
\pgfpathlineto{\pgfqpoint{1.407286in}{0.638743in}}%
\pgfpathlineto{\pgfqpoint{1.446638in}{0.635882in}}%
\pgfpathlineto{\pgfqpoint{1.492887in}{0.633369in}}%
\pgfpathlineto{\pgfqpoint{1.556801in}{0.630976in}}%
\pgfpathlineto{\pgfqpoint{1.627998in}{0.629292in}}%
\pgfpathlineto{\pgfqpoint{1.679052in}{0.628501in}}%
\pgfpathlineto{\pgfqpoint{1.726272in}{0.627978in}}%
\pgfpathlineto{\pgfqpoint{1.771947in}{0.627609in}}%
\pgfpathlineto{\pgfqpoint{1.817105in}{0.627342in}}%
\pgfpathlineto{\pgfqpoint{1.862374in}{0.627146in}}%
\pgfpathlineto{\pgfqpoint{1.908215in}{0.627000in}}%
\pgfpathlineto{\pgfqpoint{1.955016in}{0.626891in}}%
\pgfpathlineto{\pgfqpoint{2.003139in}{0.626811in}}%
\pgfpathlineto{\pgfqpoint{2.052949in}{0.626751in}}%
\pgfpathlineto{\pgfqpoint{2.104836in}{0.626707in}}%
\pgfpathlineto{\pgfqpoint{2.159239in}{0.626675in}}%
\pgfpathlineto{\pgfqpoint{2.216670in}{0.626652in}}%
\pgfpathlineto{\pgfqpoint{2.277746in}{0.626635in}}%
\pgfpathlineto{\pgfqpoint{2.343228in}{0.626624in}}%
\pgfpathlineto{\pgfqpoint{2.414092in}{0.626616in}}%
\pgfusepath{stroke}%
\end{pgfscope}%
\begin{pgfscope}%
\pgfsetrectcap%
\pgfsetmiterjoin%
\pgfsetlinewidth{0.803000pt}%
\definecolor{currentstroke}{rgb}{0.000000,0.000000,0.000000}%
\pgfsetstrokecolor{currentstroke}%
\pgfsetdash{}{0pt}%
\pgfpathmoveto{\pgfqpoint{0.564660in}{0.521603in}}%
\pgfpathlineto{\pgfqpoint{0.564660in}{2.831603in}}%
\pgfusepath{stroke}%
\end{pgfscope}%
\begin{pgfscope}%
\pgfsetrectcap%
\pgfsetmiterjoin%
\pgfsetlinewidth{0.803000pt}%
\definecolor{currentstroke}{rgb}{0.000000,0.000000,0.000000}%
\pgfsetstrokecolor{currentstroke}%
\pgfsetdash{}{0pt}%
\pgfpathmoveto{\pgfqpoint{0.564660in}{0.521603in}}%
\pgfpathlineto{\pgfqpoint{2.502160in}{0.521603in}}%
\pgfusepath{stroke}%
\end{pgfscope}%
\begin{pgfscope}%
\definecolor{textcolor}{rgb}{0.000000,0.000000,0.000000}%
\pgfsetstrokecolor{textcolor}%
\pgfsetfillcolor{textcolor}%
\pgftext[x=1.533410in,y=2.914937in,,base]{\color{textcolor}{\sffamily\fontsize{12.000000}{14.400000}\selectfont\catcode`\^=\active\def^{\ifmmode\sp\else\^{}\fi}\catcode`\%=\active\def%{\%}$n=10^6$}}%
\end{pgfscope}%
\begin{pgfscope}%
\pgfsetrectcap%
\pgfsetroundjoin%
\pgfsetlinewidth{1.505625pt}%
\definecolor{currentstroke}{rgb}{0.000000,0.000000,0.000000}%
\pgfsetstrokecolor{currentstroke}%
\pgfsetdash{}{0pt}%
\pgfpathmoveto{\pgfqpoint{1.836489in}{2.649691in}}%
\pgfpathlineto{\pgfqpoint{1.975378in}{2.649691in}}%
\pgfpathlineto{\pgfqpoint{2.114267in}{2.649691in}}%
\pgfusepath{stroke}%
\end{pgfscope}%
\begin{pgfscope}%
\definecolor{textcolor}{rgb}{0.000000,0.000000,0.000000}%
\pgfsetstrokecolor{textcolor}%
\pgfsetfillcolor{textcolor}%
\pgftext[x=2.225378in,y=2.601080in,left,base]{\color{textcolor}{\sffamily\fontsize{10.000000}{12.000000}\selectfont\catcode`\^=\active\def^{\ifmmode\sp\else\^{}\fi}\catcode`\%=\active\def%{\%}$\pi_u$}}%
\end{pgfscope}%
\begin{pgfscope}%
\pgfsetbuttcap%
\pgfsetroundjoin%
\pgfsetlinewidth{1.505625pt}%
\definecolor{currentstroke}{rgb}{0.000000,0.000000,0.000000}%
\pgfsetstrokecolor{currentstroke}%
\pgfsetdash{{5.550000pt}{2.400000pt}}{0.000000pt}%
\pgfpathmoveto{\pgfqpoint{1.836489in}{2.445834in}}%
\pgfpathlineto{\pgfqpoint{1.975378in}{2.445834in}}%
\pgfpathlineto{\pgfqpoint{2.114267in}{2.445834in}}%
\pgfusepath{stroke}%
\end{pgfscope}%
\begin{pgfscope}%
\definecolor{textcolor}{rgb}{0.000000,0.000000,0.000000}%
\pgfsetstrokecolor{textcolor}%
\pgfsetfillcolor{textcolor}%
\pgftext[x=2.225378in,y=2.397223in,left,base]{\color{textcolor}{\sffamily\fontsize{10.000000}{12.000000}\selectfont\catcode`\^=\active\def^{\ifmmode\sp\else\^{}\fi}\catcode`\%=\active\def%{\%}$\pi_l$}}%
\end{pgfscope}%
\end{pgfpicture}%
\makeatother%
\endgroup%

   \input{RBAFIG/RBA n=10^12.pgf}
\end{center}
\subsection{Choosing Degree of Prior Certitude}
This is a very interesting question. Considering the S\&P 500 for our discussion, 30 years is a long time with the dataset containing 7559 
trade days. An initial thought is to first start with the baseline prior and then calculate one's posterior having 
$d=\lambda\E[\sigma]=\lambda\sqrt{n}\E[\sqrt{p_s(1\!\m\!p_s)}]$. However, by doing this ones is changing the hyperparameter, $d$, based on the data 
and cannot define $d$ before the data is collected.
What one can do instead is compute the $\lambda\sqrt{n}\E[\sqrt{p_s(1\!\m\!p_s)}]$ of their baseline prior. 
For the S\&P 500 analysis $d=100$, $\lambda\approx2.93$ and for $d=30$ is $\lambda\approx0.88$. 
However, this again is not the point of robust Bayesian analysis since 
the amount of reasonable priors do not increase proportionally to the $\sqrt{n}$. 
Moreover, the point of robust Bayesian statistics is not to be a worse way of computing confidence intervals. Using 
confidence intervals in robust Bayesian statistics allows at least $\gamma$ credibility\footnote{$\text{Crl}_{\gamma}(\theta;\pi\in\Pi)\in\mathcal{R}(\text{Crl}_{\gamma}(\theta))$} that the 
parameter of interest, $\theta$ is in \mbox{$\mathcal{R}(\text{Crl}_{\gamma}(\theta))$}. 
So when using credible intervals scaling $d$ by $\sqrt{n}$ is unneeded and just allows one's 
credible intervals for their parameters to be excessively vague.
\subsection{$[x_f,x_s]$ Region}
From the Analysis of the S\&P 500 we see that it gets a flat top. 
$\pi_u$ for $x\in[x_f,x_s]$ becomes flat after around 50 to 100 datapoints given one gets both success and failures.
This means that for large $n$ only 2 or 3 points are actually sampled and since the region is bounded between \mbox{Mode $B_f$}.
Moreover, the length of $[x_f,x_s]$ is less than $\frac{d}{n+d}$ while $\sigma\propto\frac{1}{\sqrt{n}}$ so the interesting 
transition region of $\pi_u$ from $B_f$ to $B_s$ becomes negligible as $n$ becomes large. 
\subsection{Conclusion}
This paper has derived the basic descriptive statistics for robust Bayesian analysis of the beta distribution and computed all the steps to create 
an very effective graphing algorithm for $\pi_l$ and $\pi_u$. Additionally, it showed how robust Bayesian analysis can
be used in the real world by determining the likelihood that  the S\&P 500 will be higher at the end of the next trading day than it was at the last trading day is between 52\%$-$55\%.
\bibliographystyle{plainnat}
\bibliography{RBA.bib}
\end{document}